% ok auch von Werner, 29.9., 15:46
% update nach Absprache mit Werner , Manuel (blue)
% revision, Manuel (blue)
% revision, Werner (green)
% revision, Reinhard (red)
%
% final polishing, submission version Manuel 16 June 2014 after acceptance of [3] in LAA :-))
% very minor update by Werner 13. Juni 2014
% major update Manuel 12 June 2014
% minor update by Werner 6. Juni 2014
% updated by Manuel 3 Juni 2014
% intro by Reinhard 2 Juni 2014
%
% minor update Manuel 20 Maerz 2014
% major update Werner 19 Maerz 2014
% major update Manuel 26 Feber 2014
% update Werner 17 und 25 Feber 2014
% slight update Manuel 2 Feber 2014
%
%\documentclass[12pt]{article}
\documentclass[reqno]{amsart}
%\usepackage{a4,amssymb,amsmath,color,graphicx}%,scalerel
\usepackage{latexsym,amsmath,amssymb,graphicx,slashbox}
\usepackage{color,url,multirow}
\usepackage{epstopdf}


\newcommand{\va}{\left(\begin{array}{c}}
\newcommand{\ve}{\end{array}\right)}
\newcommand{\be}{\begin{eqnarray}}
\newcommand{\ee}{\end{eqnarray}}
\newcommand{\di}{\textup{diag}}
%\newcommand\scale[2]{\vstretch{#1}{\hstretch{#1}{#2}}}
\def\beq#1{\begin{equation}\label{#1}}
\def\eeq{\end{equation}}
%\def\rad{{}\kern-.2em\otimes\kern-.97em\oplus}
\def\rad{{}\kern.5em\makebox[0pt]{$\otimes$}\makebox[0pt]{$\oplus$}\kern.5em{}}
\def\srad{{\scriptstyle{{}\kern.4em\makebox[0pt]{$\scriptstyle\otimes$}\makebox[0pt]{$\scriptstyle\oplus$}\kern.4em{}}}}
\def\ssrad{{\scriptstyle{{}\kern.4em\makebox[0pt]{$\scriptscriptstyle\otimes$}\makebox[0pt]{$\scriptscriptstyle\oplus$}\kern.4em{}}}}
%\scale{.6}
\def\R{{\mathbb R}}
\def\N{{\mathbb N}}
\def\a{\mathbf a}
\def\b{\mathbf b}
\def\c{\mathbf c}
\def\e{\mathbf e}
\def\i{\mathbf i}
\def\u{\mathbf u}
\def\v{\mathbf v}
\def\w{\mathbf w}
\def\x{\mathbf x}
\def\y{\mathbf y}
\def\z{\mathbf z}
\def\o{\mathbf o}
\def\Ab{{\mathsf A}}
\def\Bb{{\mathsf B}}
\def\Cb{{\mathsf C}}
\def\Db{{\mathsf D}}
\def\Eb{{\mathsf E}}
\def\Fb{{\mathsf F}}
\def\Gb{{\mathsf G}}
\def\Hb{{\mathsf H}}
\def\Ib{{\mathsf I}}
\def\Jb{{\mathsf J}}
\def\Kb{{\mathsf K}}
\def\Lb{{\mathsf L}}
\def\Mb{{\mathsf M}}
\def\Nb{{\mathsf N}}
\def\Ob{{\mathsf O}}
\def\Pb{{\mathsf P}}
\def\Qb{{\mathsf Q}}
\def\Rb{{\mathsf R}}
\def\Sb{{\mathsf S}}
\def\Tb{{\mathsf T}}
\def\Ub{{\mathsf U}}
\def\Vb{{\mathsf V}}
\def\Wb{{\mathsf W}}
\def\Xb{{\mathsf X}}
\def\Yb{{\mathsf Y}}
\def\Zb{{\mathsf Z}}
\def\reff#1{(\ref{#1})}
\def\rank{\mbox{rank\,}}
\def\ran2{\mbox{rank}^{(2)}\,}
\def\T{^\top}
\def\cpr{\mbox{cpr }}
\def\bea#1{\begin{array}{#1}}
\def\ea{\end{array}}\def\ignore#1{}
\parindent0ex
\sloppy
\renewcommand{\arraystretch}{1.1}
\renewcommand{\baselinestretch}{1.5}
\newtheorem{defn}{Definition}[section]
\newtheorem{thm}{Theorem}[section]
\newtheorem{lem}{Lemma}[section]
\newtheorem{exa}{Example}[section]
\newtheorem{cly}{Corollary}[section]
\newtheorem{rem}{Remark}[section]
\newtheorem{prop}{Proposition}[section]
\newtheorem{obs}{Observation}[section]
\newtheorem{alg}{Algorithm}[section]
%\newenvironment{proof}{\noindent {\bf Proof. }}{\hfill $\Box$ \\ \medskip}
\newcommand{\bep}{\begin{proof}}
\newcommand{\ep}{\end{proof}}
\def\feta{\pmb{\eta}}
\def\lk{\left\{}
\def\rk{\right \} }
\definecolor{dgreen}{rgb}{0,.7,0}
\definecolor{purple}{rgb}{0.5,0,0.5}
\def\cor#1{{#1}} %IMB
\def\corr#1{{#1}}

\def\cog#1{{#1}} %Werner Schachinger
\def\cogg#1{{#1}}
%\def\cogg#1{{\color{dgreen}#1}} %Werner Schachinger
\def\cop#1{{#1}} %Reinhard Ullrich

\def\ru_cor_1#1{{#1}} %Reinhard Ullrich Korrekturen 1. Revision \def\ru_cor_1#1{{\color{red}#1}}
\def\ws_cor_1#1{{#1}} %WS Korrekturen 1. Revision \def\ws_cor_1#1{{\color{dgreen}#1}}
\def\corrr#1{{#1}} %IMB 1. Revision \def\corrr#1{{\color{blue}#1}} %IMB 1. Revision

\def\qf#1#2{#1\T  #2 #1}
\def\q{\mathbf q}
\def\supp#1{I(#1)}
\def\csn{{{\mathcal S}^n}}
%\def\co {\mathcal C}
%\def\cm{{\mathcal C}^*}
\def\co{\mathcal{COP}}
\def\cm{\mathcal{CP}}

\long\def\/*#1*/{}


\begin{document}

\title{New lower bounds and asymptotics for the cp-rank}

\author[I. M. Bomze]{Immanuel M. Bomze$^\star$}
\thanks{%Dept. of Statistics and Decision Support Systems,
${}^\star$ University of Vienna, Austria}

\author[W. Schachinger]{Werner Schachinger$^\star$}

\author[R. Ullrich]{Reinhard Ullrich$^\star$}


\begin{abstract}
\normalsize
Let $p_n$ denote the largest possible cp-rank of an $n\times n$ completely positive matrix. {This matrix parameter has its significance both in theory and applications, as it sheds light on the geometry and structure of the solution set of hard optimization problems in their completely positive formulation.}
Known bounds for $p_n$ are $s_n=\binom{n+1}2-4$, the current best upper bound, and the Drew-Johnson-Loewy (DJL) lower bound $d_n=\big\lfloor\frac{n^2}4\big\rfloor$. The famous DJL conjecture (1994) states that $p_n=d_n$. Here we show $\corr{p_n=\frac {n^2}2 +{\mathcal O}\big(n^{3/2}\big) = 2d_n+{\mathcal O}\big(n^{3/2}\big)}$, and construct counterexamples to the DJL conjecture for all $n\ge {12}$ (for orders seven through eleven counterexamples were already given in \cite{BSU}).
 \end{abstract}

\subjclass{Primary: 15B48; Secondary: 90C25, 15A23}
\keywords{copositive optimization, completely positive matrices, nonnegative factorization}%

\date{\small This version: \today}

%\begin{keywords} copositive optimization, completely positive matrices, cp-rank, nonnegative factorization\end{keywords}

%\begin{AMS} , 90C25, 15A23 \end{AMS}

\maketitle

%\begin{document}
\parindent0ex
%\pagestyle{empty}

\section{{Introduction: motivation, notations}}

\subsection{{Motivation: The cp-cone and copositive optimization}}

In this article we consider completely positive matrices $\Mb $ and their cp-rank.
An $n\times n$ matrix $\Mb $ is said to be {\em completely positive} if there exists a nonnegative  (not necessarily square) matrix $\Vb$ such that $\Mb =\Vb\Vb\T$ {($\hbox{}\T$ denotes transposition)}. \corr{This form of factorization is more restrictive than the general question of \corrr{(exact) nonnegative matrix factorization mainly used in multivariate data analysis and machine learning~\cite{Vava09a},} where $\Mb$ could be rectangular \corrr{and the question is to find a factorization $\Mb = \Wb\Hb$ with rectangular matrices $\Wb$ and $\Hb$ of different sizes with no negative entries under some (e.g., rank) restrictions; see also, e.g.,~\cite{Berr07a} for a widely popular approximate version of this concept.} Here, complete positivity implies \corrr{symmetry, positive-semidefiniteness and of course also} nonnegativity of all entries.}
Typically, a completely positive matrix \corr{$\Mb =\Vb\Vb\T$} may have many such factorizations, and
the {\em cp-rank} of $\Mb $, $\cpr \Mb $, is the minimum number of columns in such a nonnegative factor $\Vb$. {The set $\cm _n$ of all completely positive $n\times n$ matrices  forms a proper cone (i.e., it is pointed, convex, and solid in the sense that it has nonempty interior). With respect to the {\em Frobenius} inner product $\langle \Ab, \Bb\rangle := \mathrm{trace}(\Ab\Bb\T)$, this cone is dual to the cone  $\co _n$} of
{symmetric} {\em copositive matrices} {of order $n$}. An $n\times n$ matrix $\Sb$ is said to be copositive if $\qf{\x} \Sb \ge 0$ for every nonnegative vector $\x\in \R^n$.

Copositive and completely positive matrices are central in the rapidly evolving field of {\em copositive optimization} which links discrete and continuous optimization, and has
numerous real-world applications. For recent surveys and structured bibliographies, we refer to~\cite{Bomz11a,Bomz11b,Bure11,Duer10}, and for a fundamental text book to~\cite{Berm03}.

A {conic optimization} problem of the form
\beq{pp} \inf \{ \langle\Cb ,\Xb\rangle : \langle \Ab_i ,\Xb\rangle = b_i ,\,  i \in \lk 1,\ldots, m\rk , \, \Xb \in \cm _n \} \eeq
{is called {\em completely positive optimization problem}, but sometimes also {\em copositive optimization problem}, because} the corresponding dual problem is given as
\beq{dp} \sup \left\{ \sum_{i=1}^m b_i y_i : \y \in \R^m ,\, \Sb = \Cb - \sum_{i=1}^m y_i\Ab_i \in \co _n\right\} .\eeq
{Both problems consist in optimizing a linear form over a feasible set which can be described as the intersection of an affine subspace with one of the cones $\co _n$ or $\cm_ n$. Hence at least one optimal solution (if this exists at all) must be contained in the boundary of these cones \ws_cor_1{(unless the affine subspace is a singleton)}. Moreover, if strong duality for~\reff{pp} and~\reff{dp} holds, then there exists a primal-dual optimal pair $(\Xb^*, \Sb^*)\in \cm _n\times \co _n$ with $\langle\Sb^*, \Xb^*\rangle = 0$ or $\Sb^*\perp\Xb^*$, which relation can be exploited to obtain information about $\Xb^*$ if we have some knowledge on $\Sb^*$.}

{As remarked above, the conic primal-dual pair~\reff{pp} and~\reff{dp} serves to reformulate NP-hard optimization problems. Since everything else is linear, it is quite obvious that this approach shifts the whole complexity of the hard optimization problem into the (boundaries of the) cones $\cm_n$ and $\co _n$. These boundaries are much more complex than the boundaries of the symmetric, self-dual cones used in polynomial-time conic optimization (such as Linear or Semidefinite Optimization, or optimization over the Minkowski cone). For instance, while the boundary of the semidefinite cone
consists of matrices which are rank-deficient, the boundary of the completely positive cone $\cm_n$ also contains nonsingular matrices like the identity matrix, or matrices with all entries strictly positive like the all-ones matrix. So, neither linear constraints on the entries nor rank restrictions are sufficient to characterize or elucidate geometric properties of completely positive matrices. Therefore, the cp-rank was early recognized as a useful matrix parameter} to shed more light upon the structure and the properties of {completely positive matrices, and consequently has received considerable attention by researchers over the past decades.}

Determining the maximum possible cp-rank of $n\times n$ completely positive matrices,
$$p_n := \max\lk \cpr \Mb  : \Mb \in \cm _n \rk ,$$
is still an open problem for general $n$. It is known \cite[Theorem~3.3]{Berm03} that
$p_n=n$ if $n\le 4$,  whereas this equality does no longer hold for $n\ge 5$.
Let $d_n := \big\lfloor \frac{n^2}4\big\rfloor$ and $s_n := {n+1\choose 2}-4$. For $n\ge 5$, it is known that\ws_cor_1{~\cite{Shak13b}}
\beq{pnbracket} d_n\le p_n \le s_n \, ,\eeq
and that $d_n=p_n$ in case $n=5$~\cite{Shak13a}.
It is still unknown whether $d_6=p_6$\ws_cor_1{.}
% although \ru_cor_1{the inequality in~\reff{pnbracket} was improved} in the recent paper~\cite{Shak13b} where also the upper bound %$p_n\le s_n$ was established for the first time.


The famous Drew-Johnson-Loewy (DJL) conjecture~\cite{Drew94} is by now twenty years old. It states that $d_n=p_n$ is true for all $n\ge 5$, and some evidence in support of the DJL conjecture is found in~\cite{Berm98,Drew96a,Drew94,Loew03}, see also \cite[Section 3.3]{Berm03}. In {a} recent paper~\cite{BSU} it {was} shown that the DJL conjecture does not hold for orders $n$ {ranging between seven and eleven} by constructing examples which establish $p_n > d_n$.


\subsection{{Notations, terminology and paper structure}}

Some notation and terminology:
we abbreviate $[r\!:\!s]=\lk r,r+1,\ldots ,s\rk$ for integers $r\le s$.
Let $\e_i\in \R^n$ be the $i$th column vector of the $n\times n$ identity matrix $\Ib_n$ and ${\feta_n}=\sum_{i=1}^n \e_i$.
By ${\Eb_n}={\feta_n\feta_n\T}$ we denote the $n\times n$ matrix of all ones. The nonnegative orthant is denoted by $\R^n_+$ which contains the standard simplex
$$\Delta_n :=\lk \x\in \R^n_+ : \feta_n\T\x=1\rk .$$
The matrix $\textup{Diag}(\y)$ is a diagonal matrix {containing} the entries of $\y$ on the diagonal.
The Kronecker product is denoted by $\otimes$, and
$$\Ab \oplus \Bb = \left [ \bea{cc}\Ab &\Ob\\ \Ob\T &\Bb\ea \right ]$$
means the direct sum of two square matrices. \ru_cor_1{The letter $\Ob$ denotes the zero matrix of appropriate order.} {For a given $\x\in \R^k_+$, we define the {\em zero-norm} $\|\x\|_0$ as the number of positive entries $x_i>0$.}
Given a square matrix $\Sb\in  \co _n$, we will use the phrase ``zero(es) of $\Sb $" as an abbreviation of ``zero(es) $\x\in \Delta_n$ of the quadratic form $\x\T\Sb\x$"; this terminology differs slightly from that in~\cite{Hild14a} {but is more convenient for our purposes}.

The paper is organized as follows: In Section~2 we look at copositive matrices $\Sb $ with finitely many (but many) {zeroes}. Such matrices $\Sb $ lie on the boundary of the copositive cone, and elementary conic duality therefore tells us that there are nontrivial completely positive matrices $\Mb $ such that $\Mb \perp \Sb $.  There is a strong connection between the zeroes of $\Sb$ and the cp-rank of $\Mb$, which is established through Lemma~\ref{lowerb}. Lemma~\ref{ran2rec} deals with cp-ranks of \corr{Khatri-Rao-like products (defined in Subsection~\ref{KRprod})} of matrices, which are necessary to make assertions about cp-ranks of high-{order} matrices. {Combination of these auxiliary results} culminates in Theorem~\ref{ninequarters} {and in Corollary \ref{clyMain}, which refutes the DJL conjecture} for $n\ge 7$ and shows that the largest possible cp-rank $p_n$ lies asymptotically much closer to the upper bound $s_n$ than to the lower bound $d_n$.

Section~3 improves the lower bound for $p_n$ in the following way: in Section~2 only identity matrices \cogg{are} used as building blocks to construct matrices of higher {order}. This is sufficient to prove the assertions of Section 2, but better results can be obtained by using, {as building blocks,} matrices with \cogg{cp-ranks that exceed their orders.}
%higher cp-ranks than the identity matrices.
Some of these building-block-matrices are new in the literature, some of them were already used in~\cite{BSU}.  To further illustrate the advantage of the approach in this article, an explicit construction of a matrix of order twelve with high cp-rank is presented in an appendix.
\corr{Note that in contrast to~\cite{BSU}, for general order $n$, we need not construct the matrices explicitly but rather can invoke the existence result in Lemma~\ref{lowerb}}.

\newpage

\section{Main results}
\subsection{{Rank, two-rank and cp-rank}}

\cogg{Our method of finding matrices of high cp-rank builds upon two observations:
(1) For certain matrices $\Mb\in\cm_n$, only multiples of vectors from a finite set $\{\u_1,\ldots,\u_k\}$ may appear as columns of a factor $\Vb$ in any factorization $\Mb=\Vb\Vb\T=\sum_{i=1}^ky_i\u_i\u_i\T$, where $\y=[y_1,\ldots,y_k]\T\in\mathbb{R}^k_+$. This property is shared by all matrices in a certain convex subcone of $\cm_n$ determined by the set $\{\u_1,\ldots,\u_k\}$.
(2) \corr{This} subcone contains matrices \corr{with} cp-rank \corr{bounded below} by a number computable from the set $\{\u_1,\ldots,\u_k\}$;
\corr{so} the cp-rank will be high if \corr{this} number is large.\\
\corr{This argument} is made precise in the following \corr{results,} sta\corr{r}ting \corr{from the observation} in a more general context, that
a convex cone spanned by some finite set of vectors of rank $r$ always contains a vector \corr{which} is {\em not} a positive linear combination of {\em less than} $r$ vectors from \corr{this} finite set; \corr{a converse of Caratheodory's theorem in some sense}.}

\begin{lem}\label{rankl} Let $V$ be a real vector space, let $\{\v_i:i\in[1\!:\! k]\}\subseteq V$ be a set of vectors of rank $r$, and define for $\y\in\mathbb{R}^k_+$
$$P_\y:=\left\{\x\in\mathbb{R}^k_+:\sum_{i=1}^kx_i\v_i=\sum_{i=1}^ky_i\v_i \right\}.$$
Then there exists $\y\in\mathbb{R}^k_+$ such that
$$\min_{\x\in P_\y}\|\x\|_0\, =\, r\, .$$
\end{lem}
\bep
First we show that $\min_{\x\in P_\y}\|\x\|_0 \le r $ for all $\y\in \R^n_+$ (this is basically Caratheodory's theorem, we include the short argument for the readers' convenience). To this end, choose an $\x\in P_\y$ with $m=\|\x\|_0$ minimal over $P_\y$. We assume without loss of generality $x_i>0$ for all $i\le m$.
If $m>r$ would hold, then there were $\mu_i\in \R$ with $\sum\limits_{i=1}^m \mu_i\v_i =\o$ with some $\mu_i > 0$. Further without loss of generality we assume {(for some $s\in [1\!:\!m]$) that $\mu_i\le 0$ for $i<s$ while $\mu_i > 0$} and $\frac {x_i}{\mu_i} \ge \frac {x_m}{\mu_m} > 0$ for all $i\in [s\!:\!m]$. Define $z_i := x_i - \frac {x_m}{\mu_m}\mu_i \ge0$ for all $i\in[1\!:\!m]$ and $z_i:=0$ for $i>m$, so that $\|\z\|_0 \le m-1$ (as also $z_m=0)$. But straightforward calculations show $\sum_i z_i\v_i = \sum_i x_i\v_i$, so $\z\in P_\y$, in contradiction to the assumptions.
Next we use the fact that a vector space {over an} infinite scalar field is never the union of a finite number of proper subspaces, see~\cite[p.211]{H95}. %\ep
Define the cone $C:=\left\{\sum_{i=1}^m y_i\v_i:\y\in\mathbb{R}^m_+\right\}\subseteq V$ and observe  that the linear subspace $L=C-C$ is $r$-dimensional.
If we had $\min\limits_{\x\in P_\y}\|\x\|_0<r$ for all $\y\in \mathbb{R}^m_+$, then $C$ (and thus also $L$) would have to be a subset of
$$U:=\bigcup_{\substack{I\subseteq[1 : m] \\ |I|\le r-1}}
\left\{\sum_{i=1}^m x_i\v_i:\x\in\mathbb{R}^m , \, {x_i=0 \mbox{ for all }i\in [1\!:\!m]\setminus I}\right\},$$
which is impossible, since $U$ is a union of finitely many proper subspaces of $L$ (of dimension at most $r-1$).
\ep

For a matrix $\Ab=[\a_1,\ldots,\a_k]\T$ we let $\Ab^{\langle 2\rangle}:=[\a_1\otimes\a_1,\ldots,\a_k\otimes\a_k]\T$, and {define the {\em two-rank} of $\Ab$ as}
$$\ran2 \Ab:=\rank \Ab^{\langle 2\rangle}\, .$$
For illustration, denote by $\Bb_i = \e_i\e_i\T\in \R^{n\times n}$. Then $\Ib_n^{\langle 2\rangle} = [\Bb_1 | \cdots | \Bb_n]$.
Note that always $\ran2 \Ab\ge\rank \Ab$ with equality if $\rank \Ab = k$, i.e., if $\Ab$ itself has full row rank, then also $\Ab^{\langle 2\rangle}$ has (the same) full row rank. Furthermore we note for later use the trivial relations $\ran2 (\alpha \Ab) = \ran2 \Ab$ if $\alpha>0$,
$$ \ran2 %\left [ \bea{c}\Ab \\ \Bb \ea \right ]
\bigg[   { \Ab \atop \Bb }\bigg]
\ge \ran2 \Bb\, ,$$
and a slightly less trivial one: $\ran2 [\Ab | \Bb ]\ge \ran2 \Bb$.
\begin{lem}\label{lowerb}
Let $\Ub=[\u_1,\ldots,\u_k]\T\in\mathbb{R}^{k\times n}_+$, where $\{\u_1,\ldots,\u_k\}$ are all the zeroes of some copositive matrix $\Sb\in\co _n$.\\
Then there \corr{exists}
a diagonal matrix $\Db=\textup{Diag}(\y)$ with $\y\in\mathbb{R}^k_+$ such that the completely positive matrix $\Mb=\Ub\T\Db\Ub$ satisfies $\cpr \Mb = \ran2 \Ub$.
\end{lem}
\bep


{Consider any $\Mb =\Ub\T\textup{Diag}(\y)\Ub$.} We observe that $\langle \Mb,\Sb\rangle =0$, i.e., that $\Mb\perp\Sb$  holds. Therefore by \cite[Lemma~2.1]{BSU} we conclude that any cp-factorization of $\Mb$ is of the form
$$\Mb=\Ub\T\textup{Diag}(\x)\Ub=\sum_{i=1}^kx_i\u_i\u_i\T$$ with some $\x\in\mathbb{R}^k_+$.
\cogg{For any $\x$ corresponding to a minimal cp-factorization of $\Mb$ we then have $\cpr\Mb=\|\x\|_0$.}
Since the rank of the set $\{\u_i\u_i\T:i\in[1\!:\!k]\}$ equals $\ran2\Ub$, as is seen by identifying $\u_i\u_i\T$ with $\u_i\otimes\u_i={\rm vec}(\u_i\u_i\T)$, we can invoke Lemma~\ref{rankl} to obtain the desired conclusion.
\ep

\subsection{{Direct sums and Khatri-Rao-like products}}\label{KRprod}


For matrices $\Ub\in\mathbb{R}^{k\times n}$ and $\Vb\in\mathbb{R}^{\ell\times m}$ we
construct the following $k\ell\times (n+m)$-matrix,
denoted as
$\Ub\rad \Vb=\left[\Ub\otimes\pmb{\eta}_\ell |  \pmb{\eta}_k\otimes \Vb \right]$; \corr{recall that} $\pmb{\eta}_d$ denotes the all ones vector in $\mathbb{R}^d$.
Note that both $\Ub\otimes\Vb$ and $\Ub^{\langle 2\rangle}\rad\Vb^{\langle 2\rangle}$ are, up to permutations of columns, submatrices of $(\Ub\rad \Vb)^{\langle 2\rangle}$, and all these matrices have the same number $k\ell$ of rows. Further note that using the Khatri-Rao product~$\star$, see. e.g.~\cite{KhatriRaoRef},  we can write
$$\Ub\rad \Vb =[\Ub |\pmb{\eta}_k]\star [\pmb{\eta}_\ell | \Vb]\quad\mbox{and}\quad (\Ub^{\langle 2\rangle})\T= [\u_1 | \u_2 | \cdots | \u_k] \star [\u_1 | \u_2 | \cdots | \u_k] \, .$$
Recall that a matrix $\Ab\in\mathbb{R}^{k\times n}_{+}$ is {\em row-stochastic} if $\Ab\pmb{\eta}_n= \pmb{\eta}_k$ holds.

\begin{lem}\label{ran2rec}
Let $\alpha>0$ and $\beta >0$ and consider two row-stochastic matrices $\Ub\in\mathbb{R}^{k\times n}$ and  $\Vb\in\mathbb{R}^{\ell\times m}$.
Then for the $k\ell\times (n+m)$-matrix
$\Wb=(\alpha\Ub)\rad (\beta\Vb)$
and the $(k+\ell)\times (n+m)$-matrix $\widetilde \Wb= \Ub\oplus \Vb$
 we have
\begin{enumerate}
  \item [(a)] $\rank \Wb=\rank \Ub+\rank \Vb-1$ and $\frac 1{\alpha +\beta}\Wb$ is row-stochastic,
  \item [(b)] $\ran2 \Wb\ge \rank \Ub\,\rank \Vb+\ran2 \Ub -\rank \Ub+\ran2 \Vb-\rank \Vb$,
  %$\ran2 \Wb\ge (\rank \Ub-1)(\rank \Vb-1)+\ran2 \Ub +\ran2 \Vb -1$
  \item [(c)] $\rank \widetilde \Wb=\rank \Ub+\rank \Vb,~\ran2 \widetilde \Wb=\ran2 \Ub+\ran2 \Vb$\\ and $\widetilde \Wb$ is \cor {row-stochastic.}
  \item [(d)] If the rows of $ \Ub$ (resp.~$ \Vb$) are all the zeroes of some
    $\Sb_\Ub\in\co _n$ (resp.~$\Sb_\Vb\in\co _m$),
    then there are copositive matrices $\big\{ \Sb,\widetilde \Sb\big\}\subset\co _{n+m}$ such that
    the rows of $\frac 1{\alpha +\beta}\Wb$ are all the zeroes of $\Sb$ and
    the rows of $\widetilde \Wb$ are all the zeroes of $\widetilde \Sb$.
\end{enumerate}
\end{lem}
\bep

It is clear that $\frac 1{\alpha +\beta}\Wb$ is row-stochastic. Let $r_\Ub:=\rank \Ub $ and $r_\Vb:=\rank \Vb $. Since the rank of the first $n$ (resp.~last $m$) columns of $\Wb$ is $r_\Ub$ (resp.~$r_\Vb$), $\rank \Wb$ can be smaller than $r_\Ub+r_\Vb$ only if some nonzero linear combination of the first $n$ columns of $\Wb$ equals some linear combination of the last $m$ columns of $\Wb$.

%%% neu
So assume that there are $\x\in\mathbb{R}^n$ and
$\y\in\mathbb{R}^m$, such that
$(\Ub\otimes\pmb{\eta}_\ell)\x=(\Ub\otimes\pmb{\eta}_\ell)(\x\otimes1)=\Ub\x\otimes\pmb{\eta}_\ell$
and $(\pmb{\eta}_k\otimes \Vb)\y=(\pmb{\eta}_k\otimes \Vb)(1\otimes\y)=\pmb{\eta}_k\otimes \Vb\y$
are both equal to $\w\in\mathbb{R}^{k\ell}\setminus\{\o\}$. From $\w=\Ub\x\otimes\pmb{\eta}_\ell$ we deduce that
$w_i=w_j$ if $\lceil\tfrac i\ell\rceil=\lceil\tfrac j\ell\rceil$, and from
$\w=\pmb{\eta}_k\otimes \Vb\y$ we deduce $w_i=w_j$ if $i\equiv j \textup{ mod } \ell$, and the only
nonzero vectors satisfying both conditions are of the form $\w=c\,\pmb{\eta}_{k\ell}$ with $c\ne0$.
%%% ende neu
%Such linear combinations indeed exist, and they are all of the form $c\,\pmb{\eta}_{k\ell}$ with $c\ne0$.
Therefore $\rank \Wb =r_\Ub+r_\Vb-1$, which concludes the proof of (a).\\
Next we denote $\rho_\Ub=\ran2\Ub$ and $\rho_\Vb=\ran2 \Vb$, and assume that the rows of $\Ub$ and $\Vb$ are arranged in a way such that
the matrices $\widetilde \Ub=\Ub_{[1:r_\Ub]\times[1:n]}$, $\widetilde \Vb=\Vb_{[1:r_\Vb]\times[1:m]}$, $\widehat \Ub=\Ub_{[r_\Ub+1:\rho_\Ub]\times[1:n]}$ and $\widehat \Vb=\Vb_{[r_\Vb+1:\rho_\Vb]\times[1:m]}$ satisfy
$$\rank \widetilde \Ub =r_\Ub,\ \rank \widetilde \Vb =r_\Vb,\
\ran2 \bigg[   { \widetilde \Ub \atop    \widehat \Ub }\bigg] =\rho_\Ub,\
\ran2 \bigg[   { \widetilde \Vb \atop    \widehat \Vb }\bigg] =\rho_\Vb.$$
Moreover let $\u_1=\e_1\T \Ub$ and $\v_1=\e_1\T \Vb$ be the first rows of $\Ub$ and $\Vb$.
\renewcommand\arraystretch{0.8}
Now consider the following $(r_\Ub r_\Vb+\rho_\Ub-r_\Ub+\rho_\Vb-r_\Vb)\times(n+m)$-submatrix of
\ws_cor_1{$\Wb':=\Wb\left[
           \begin{array}{cc}
             \!\frac1\alpha\Ib_n\! &\! \Ob\! \\
             \!\Ob\! &\! \frac1\beta\Ib_m\! \\
           \end{array}
         \right]$}:
$$\overline \Wb=\left[
           \begin{array}{cc}
             \widetilde \Ub\otimes\pmb{\eta}_{r_\Vb} & \pmb{\eta}_{r_\Ub}\otimes\widetilde \Vb \\
             \u_1\otimes\pmb{\eta}_{\rho_\Vb-r_\Vb} & \widehat \Vb \\
             \widehat \Ub &\pmb{\eta}_{\rho_\Ub-r_\Ub}\otimes\v_1 \\
           \end{array}
         \right]=
         \left[
           \begin{array}{c}
             \widetilde \Ub\rad\widetilde \Vb \\
             \u_1\rad\widehat \Vb \\
             \widehat \Ub \rad\v_1 \\
           \end{array}
         \right].
$$
Noting that $\widetilde \Ub\otimes\widetilde \Vb$ is a submatrix of $(\widetilde \Ub\rad\widetilde \Vb)^{\langle 2\rangle}$, where the latter has $r_\Ub r_\Vb$ rows,
we deduce $\rank(\widetilde \Ub\rad\widetilde \Vb)^{\langle 2\rangle}= r_\Ub r_\Vb$
from
$$r_\Ub r_\Vb=\rank\widetilde \Ub\,\rank\widetilde \Vb=\rank(\widetilde \Ub\otimes\widetilde \Vb)\le\rank(\widetilde \Ub\rad\widetilde \Vb)^{\langle 2\rangle}\le r_\Ub r_\Vb.$$
Next consider the submatrix $\left[
           \begin{array}{c}
            \!\widetilde\Ub^{\langle 2\rangle}\rad\widetilde\Vb^{\langle 2\rangle}\! \\
            \!\u_1^{\langle 2\rangle}\rad\widehat\Vb^{\langle 2\rangle}\! \\
             \!\widehat\Ub^{\langle 2\rangle}\rad\v_1^{\langle 2\rangle}\! \\
           \end{array}
         \right]$
of $\overline \Wb{}^{\langle 2\rangle}\!=\!\left[
           \begin{array}{c}
            \! (\widetilde \Ub\rad\widetilde \Vb)^{\langle 2\rangle}\! \\
            \! (\u_1\rad\widehat \Vb)^{\langle 2\rangle}\! \\
            \! (\widehat \Ub \rad\v_1)^{\langle 2\rangle}\! \\
           \end{array}
         \right]$.
If for $\x\in\mathbb{R}^{r_\Ub r_\Vb}$, $\y\in\mathbb{R}^{\rho_\Vb- r_\Vb}$, $\z\in\mathbb{R}^{\rho_\Ub- r_\Ub}$ we have
$\o=[\x\T,\y\T,\z\T]\overline \Wb{}^{\langle 2\rangle}$, then also
\begin{align*}
\o={}&\x\T(\widetilde\Ub^{\langle 2\rangle}\rad\widetilde\Vb^{\langle 2\rangle})+\y\T(\u_1^{\langle 2\rangle}\rad\widehat\Vb^{\langle 2\rangle})+\z\T(\widehat\Ub^{\langle 2\rangle}\rad\v_1^{\langle 2\rangle})\\
={}&\x\T(\widetilde\Ub^{\langle 2\rangle}\rad\widetilde\Vb^{\langle 2\rangle})+(\y\T\pmb{\eta}_{\rho_\Vb-r_\Vb}\u_1^{\langle 2\rangle})\rad(\y\T\widehat\Vb^{\langle 2\rangle})
+(\z\T\widehat\Ub^{\langle 2\rangle})\rad(\z\T\pmb{\eta}_{\rho_\Ub-r_\Ub}\v_1^{\langle 2\rangle})
\end{align*}
must hold. Therefore $\y\T\widehat\Vb^{\langle 2\rangle}$ belongs to the row space of $\widetilde\Vb^{\langle 2\rangle}$, and
$\z\T\widehat\Ub^{\langle 2\rangle}$ belongs to the row space of $\widetilde\Ub^{\langle 2\rangle}$, implying $\y=\o$ and $\z=\o$, because, by assumption, the rows of both $\displaystyle \bigg[   { \widetilde \Ub^{\langle 2\rangle} \atop    \widehat \Ub^{\langle 2\rangle} }\bigg]$ and
$\displaystyle\bigg[   { \widetilde \Vb^{\langle 2\rangle} \atop    \widehat \Vb^{\langle 2\rangle} }\bigg]$ are linearly independent.
Then by linear independence of the first $r_\Ub r_\Vb$ rows of $\overline \Wb{}^{\langle 2\rangle}$ also $\x=\o$ must hold.
Thus $\ran2 \Wb\ws_cor_1{=\ran2 \Wb'}\ge\ran2\overline \Wb=r_\Ub r_\Vb+\rho_\Ub-r_\Ub+\rho_\Vb-r_\Vb$, which completes the proof of (b).\\
For the proof of (c)
%we use that diagonal equivalence does not change rank and $\ran2$. For the second assertion
we use
\cogg{that for any matrices $\Ab,\Bb$ we have $\rank (\Ab\oplus\Bb)=\rank\Ab+\rank\Bb$, and that}
the matrix $(\Ab\oplus\Bb)^{\langle 2\rangle}$ and its submatrix $\Ab^{\langle 2\rangle}\oplus\Bb^{\langle 2\rangle}$ have the same rank.
Furthermore $\widetilde \Wb\pmb{\eta}_{n+m}=\pmb{\eta}_{k+\ell}$ is easily checked.\\
Finally, for the proof of (d) we define matrices
$$
\Sb:=\left [
          \begin{array}{cc}
            \Sb_\Ub+\frac{\beta}{\alpha}\Eb_n&-\pmb{\eta}_n\pmb{\eta}_m\T  \\
            -\pmb{\eta}_m\pmb{\eta}_n\T & \Sb_\Vb+\frac{\alpha}{\beta}\Eb_m\\
          \end{array}
        \right ],\quad\textup{ and }\quad
\widetilde\Sb:=\left [
          \begin{array}{cc}
            \Sb_\Ub&\pmb{\eta}_n\pmb{\eta}_m\T  \\
            \pmb{\eta}_m\pmb{\eta}_n\T &\Sb_\Vb\\
          \end{array}
        \right ].
$$
Take any $\z=[\lambda\x\T,(1-\lambda)\y\T]\T$ with $(\x,\y)\in\Delta_n\times\Delta_m$ and $0\le\lambda\le1$. Then
$$\z\T \Sb\z=\lambda^2\x\T \Sb_\Ub\,\x+(1-\lambda)^2\y\T \Sb_\Vb\,\y
+ \frac{(\alpha+\beta)^2}{\alpha\beta}\left(\lambda-\frac{\alpha}{\alpha+\beta}\right)^2\ge 0\, ,$$
with equality if and only if $\lambda=\frac{\alpha}{\alpha+\beta}$, and $\x\T$ (resp.~$\y\T$) is one of the rows of $ \Ub$ (resp.~$ \Vb$), i.e.~if and only if $\z\T$ is one of the rows of $\frac1{\alpha+\beta} \Wb$.\\
Furthermore, with $\z$ as above, we have
$$\z\T \widetilde\Sb\z=\lambda^2\x\T \Sb_\Ub\,\x+(1-\lambda)^2\y\T \Sb_\Vb\,\y
+ 2\lambda (1-\lambda)\ge 0\, ,$$
with equality if and only if $\lambda\in\{0,1\}$, and, depending on the value of $\lambda$, either $\x\T$ is one of the rows of $\Ub$ or $\y\T$ is one of the rows of $\Vb$, i.e.~if and only if $\z\T$ is one of the rows of $\widetilde \Wb$.
\ep

\subsection{{Zeroes and characteristic triples}}

We now define the set $\mathcal{Z}$ as follows: denote by $\mathcal{R}$ all row-stochastic matrices and let
$$\mathcal{R}_0 :=\lk \Ub\in\mathcal{R} : \mbox{the rows of $\Ub$ are all the zeroes of some copositive matrix}\rk$$
as well as $\mathcal{Z} :=\lk \alpha \Ub : \alpha >0\, , \; \Ub\in\mathcal{R}_0\rk$.

The matrices in $\mathcal{Z}$ are exactly those that are needed for applications of Lemma~\ref{lowerb}.
Moreover, with Lemma~\ref{ran2rec} we have a means of constructing new elements
%\footnote{There is also another closure property which we won't use: If $\Pb$ and $\Db$ are a permutation matrix and a positive diagonal %matrix of suitable orders, then $\Ub\in\mathcal{Z}\Rightarrow\Pb\Ub\Db\in\mathcal{Z}$.}
of $\mathcal{Z}$ from known ones:
$\lk \Ub,\Vb\rk \subset\mathcal{Z}\Rightarrow\Ub\rad\Vb\in\mathcal{Z}$. For our purpose of showing the existence of matrices of large cp-rank, only certain characteristics of a matrix $\Ub\in\mathcal{Z}$ are important: (a) the number of columns of $\Ub$ (say $n$); (b) the $\rank \Ub$ (say $r$); and (c) an integer lower bound $\rho$ for $\ran2\Ub$ (where we require $\rho\ge\rank\Ub$); these three integers we collect in a \emph{characteristic triple}
$$\corr{c =(\corrr{c_1},\corrr{c_2},\corrr{c_3}) := (n,r,\rho)\, .}$$
Some $\Ub\in\mathcal{Z}$ may have more than one characteristic triple, namely if and only if $\rank \Ub < \ran2 \Ub$. By abuse of notation, \cogg{we define} a binary operation on any two characteristic triples,
\beq{tripl}(n_1,r_1,\rho_1)\rad(n_2,r_2,\rho_2):=(n_1+n_2,r_1+r_2-1,r_1r_2+\rho_1-r_1+\rho_2-r_2)\, ;\eeq
\corr{note that $1\le r_1+r_2-1\le r_1r_2+\rho_1-r_1+\rho_2-r_2$ if both $1\le r_1\le\rho_1$ and $1\le r_2\le\rho_2$ holds. The operation $\rad$ obviously} obeys the commutative and (only a little less obvious) the associative law.\footnote{{The binary operation $\rad$ on the set $\mathcal{Z}$ is associative but not commutative, but there are always permutation matrices $\Pb_1, \Pb_2$ such that $\Bb\rad\Ab=\Pb_1(\Ab\rad\Bb)\Pb_2$. Clearly, row and column permutations of $\Ub$ do neither affect $\rank \Ub$ nor $\ran2\Ub$.}}
Clearly also the binary operation
$(n_1,r_1,\rho_1)\oplus(n_2,r_2,\rho_2):=(n_1+n_2,r_1+r_2,\rho_1+\rho_2)$
is associative and commutative.
\cogg{It follows from Lemma~\ref{ran2rec} that if $c,c'$ are characteristic triples of $\Ub,\Vb\in\mathcal{Z}$, then $c\rad c'$ is a
characteristic triple of $\Ub\rad\Vb$ and $c\oplus c'$ is a characteristic triple of $\Ub\oplus\Vb$.}\\
\cogg{Our strategy is to fix a subset $\mathcal{U}\subseteq\mathcal{Z}$ together with a set $C$ of characteristic triples, containing one characteristic triple
for each $\Ub\in\mathcal{U}$, and construct the $\rad$-semigroups \corrr{$\bar{\mathcal{U}}$ and $\bar C$} generated by $\mathcal{U}$ and $C$. 
From the latter, \corr{we fix the first component, \corrr{i.e., require} $\corrr{c_1}=n$}, the column number of some \corr{$\Ub\in\corrr{\bar{\mathcal{U}}}$ accordingly picked}, and search a triple $c\in \corrr{\bar C}$ \corr{with a large} third component $\corrr{c_3} \le \ran2 \Ub$. There are no limitations on the second component $\corrr{c_2} = \rank \Ub$, and typically the chosen $\Ub$ will not have full column rank.}

We start considering semigroups generated by a single $\Ub\in\mathcal{Z}$,
and therefore define $\Ub^{\srad1}=\Ub$ and inductively $\Ub^{\srad(n+1)}=\Ub\rad \Ub^{\srad n}$ for $n\ge1$. Similarly we define $c^{\srad n}$, where $c$ is a characteristic triple.
\begin{thm}\label{t01}
Let $\Ub\in\mathcal{Z}$, and $(n,r,\rho)$ be (one of) its characteristic triple(s). Then for any
$i\in\mathbb{N}$ there is a matrix $\Mb\in\cm _{ni}$ satisfying
$$\cpr\Mb\ge\tfrac12(r-1)^2i(i-1)+(\rho-1)i+1\, .$$
\end{thm}
\bep
The result follows from
$$(n,r,\rho)^{\srad i}=(ni,(r-1)i+1,\tfrac12(r-1)^2i(i-1)+(\rho-1)i+1)\, ,$$
which is easily proved by induction, {using~\reff{tripl}.}
\ep

For any $n\ge1$ we have $\Ib_n\in\mathcal{Z}$, since
the rows of $\Ib_n$ are the only zeroes of the copositive matrix $\Eb_n-\Ib_n$. The (unique) characteristic triple of $\Ib_n$ is $(n,n,n)$.
Putting $\Ub=\Ib_n$ in Theorem~\ref{t01}, we see
$\ran2 \Ib_n^{\srad i} \ge \rho_{i,n}$ where
\beq{defroin}\rho_{i,n}:=\tfrac12(n-1)^2i(i-1)+(n-1)i+1=\tfrac{(ni)^2}{2}-ni(i+\tfrac n2)+2ni+\tfrac{i(i-3)}{2}+1\, . \eeq
Next, counterexamples to the DJL-conjecture for infinitely many $n$, and in particular for $n=12$, are presented.

\begin{exa}\label{121314}
$(p_{12}\ge37>36=d_{12})$ We have $\ran2\Ib_4^{\srad3}\ge37$ by~\reff{defroin}; more precisely, we have $(4,4,4)^{\srad3}=(12,10,37)$ and thus
there is a completely positive matrix of order 12, rank 10 and cp-rank at least $37$, which may be written as $(\Ib_4^{\srad3})\T\textup{Diag}(\x)(\Ib_4^{\srad3})$ for some $\x\in\mathbb{R}_+^{64}$. An explicit construction will be given in {Appendix A.}\\
Similarly we obtain
$p_{4i} \ge\rho_{i,4} = \tfrac9{32}(4i)^2-\tfrac38(4i)+1  > \lfloor\tfrac14(4i)^2\rfloor =d_{4i}$ for $i\ge3$ and
$p_{3n}  \ge \rho_{3,n}  =  \tfrac13(3n)^2-3n+1   >  \lfloor\tfrac14(3n)^2\rfloor  =  d_{3n}$ for $n\ge4$.
\end{exa}


We continue this argument and maximize, for fixed $N:=ni$, the second term in formula~\reff{defroin} for $\rho_{i,n}$, namely
$-ni(i+\tfrac n2)$, with the result $n^*=\sqrt{2N}, i^*=\sqrt{\frac N2}$, and
\beq{golobo}\rho_{i^*,n^*}=\cogg{\frac{N^2}2-N\sqrt{2N}+\frac94N-\frac34\sqrt{2N}+1}\, \eeq
%\beq{golobo}\rho_{i^*,n^*}=\tfrac12N^2-\sqrt{2}N^{\frac32}+\tfrac94N-\tfrac32\sqrt{\tfrac N2}+1\, \eeq
which yields good lower bounds for the cp-rank if both $i^*$ and $n^*$ are integers, i.e., if
$n=2m$ (and $N=2m^2$) for $m\in \mathbb{N}$. \corr{We will re-encounter the three leading terms of~\reff{golobo} in the estimate~\reff{satz32abs} of Corollary~\ref{clyMain} below; for an improvement see~\reff{uffabs} in Section~\ref{improbo}.}

Still better lower bounds could probably be obtained by considering products of possibly different characteristic triples (just before we only considered powers of a single characteristic triple). Let $S$ be the semigroup generated by the set of characteristic triples $\{(i,i,i):i\in\mathbb{N}\}$.
\cogg{So any $c\in S$ is a finite \corr{$\rad$-}product of \corr{these $(i,i,i)$,} allowing repetition of factors, and $\corrr{c_1-c_2}+1$ is the number of factors, counted with multiplicity.
The factorization need however not be unique, as is seen from the example $(12,10,30)=(1,1,1)\rad(5,5,5)\rad(6,6,6)=(2,2,2)\rad(3,3,3)\rad(7,7,7)$.}
The best lower bound for $p_n$ that we can get from $S$ is then
\beq{rhoundb} b_n:=\max \lk \corrr{c_3 : c_1=n}\, ,\,\ c\in S\rk \, .\eeq
\begin{lem}\label{twofactors} The maximum $b_n$ is for some $j\ge1$ attained at a characteristic triple $c$ of the form
$c=(i_1,i_1,i_1)\rad\cdots\rad(i_j,i_j,i_j),$
where $i_1\le\cdots\le i_j$ and $i_j-i_1\le1$.
\end{lem}
\bep
There is nothing to show if $j=1$. \corrr{So assume $j\ge 2$ and $i_j-i_1>1$. We use the following property holding for any three characteristic triples $c,c',c''$:}
\beq{monoprop}
\corrr{\mbox{if}\quad c_1=c_1'\, , \; c_2=c_2' \quad\mbox{and}\quad c_3 < c_3'\, ,\quad\mbox{ then }\quad
(c\rad c'')_3 <(c'\rad c'')_3\, .}\eeq
\corrr{By assumption that $j\ge 2$ with $i_j-i_1>1$, the maximizing triple can be written % we have the representation
as $c=(i_1,i_1,i_1)\rad(i_j,i_j,i_j)\rad c''$ for some characteristic triple $c''$ (which is not there if $j=2$).
We compare $c$} with
$\tilde c:=(i_1+1,i_1+1,i_1+1)\rad(i_j-1,i_j-1,i_j-1)\rad c''$. \corrr{Applying~\reff{tripl} we obtain $c_1 = \tilde c_1$ as well as $c_2 = \tilde c_2$, and thus, by~\reff{monoprop},
$c_3 < \tilde c_3$, which contrasts} maximality of \corrr{$c_3$}. \ep

\cogg{We remark that the characteristic triples that maximize \eqref{rhoundb} are in general not uniquely determined by Lemma~\ref{twofactors}. For example, $b_{37}=442$ is attained twice, at $(7,7,7)^{\srad 3}\rad(8,8,8)^{\srad 2}\!=\!(37,33,442)$ and at $(9,9,9)^{\srad 3}\rad(10,10,10)\!=\!(37,34,442)$.}

\subsection{{New bounds for the cp-rank}}

In the following theorem we provide precise asymptotic estimates for $b_n$ as defined in \eqref{rhoundb}.

\begin{thm}\label{ninequarters}
For $n\ge5$, we have
\beq{asugl}\cogg{-\sqrt{2n}+\corr{\frac1{16}}
\le b_n-\frac{n^2}2+ n\sqrt{2n}-\frac94n\le -\frac58\sqrt{2n}+\frac32\, .}\eeq
%\beq{asugl}-2\sqrt{\tfrac n2}+\corr{\tfrac1{16}}
%\le b_n-\tfrac12{n^2}+ \sqrt{2}\, n^{\frac32}-\tfrac94n\le -\tfrac54\sqrt{\tfrac n2}+\tfrac32\, .\eeq
Moreover $b_n\le\cpr \Mb$ for some $\Mb\in\cm _n$ of the form
$\Mb=\Ub\T\Db\Ub,$
 where $\Db$ is a non\-negative diagonal matrix and $\Ub\in\mathcal{Z}$ is a binary matrix, i.e., has all entries in $\lk 0,1\rk$.
\end{thm}
\bep
From Lemma~\ref{twofactors} we know that \corr{$b_n=\corrr{c_3}$ for} some characteristic triple \corr{$c\in S$} of the form
$$c= (n,n-k+1,\rho_{m,k,i})=(m,m,m)^{\srad i}\rad(m+1,m+1,m+1)^{\srad k-i},$$
where $m\ge 1$, $k\ge1$, $1\le i\le k$ and $n=mi+(m+1)(k-i)=mk+k-i$.
Then the binary matrix $\Ub:=\Ib_m^{\srad i}\rad\Ib_{m+1}^{\srad k-i}\in\mathcal{Z}$ satisfies
$\ran2\Ub\ge b_n$, and by Lemma~\ref{lowerb} there is a nonnegative diagonal matrix $\Db$ such that we have $b_n\le\cpr \Ub\T\Db\Ub$, which settles the second assertion of the theorem. We now turn to the asserted inequalities.
Putting $r_1=(m-1)i+1$, $r_2=m(k-i)+1$, $\rho_1=\rho_{i,m}$ and $\rho_2=\rho_{k-i,m+1}$ from~\reff{defroin}
yields
\begin{align*}
\rho_{m,k,i}={}&r_1r_2+ \rho_1-r_1+\rho_2-r_2\\
={}&(i(m\!-\!1)\!+\!1)((k\!-\!i)m\!+\!1)+\tfrac12(m\!-\!1)^2i(i\!-\!1)+\tfrac12m^2(k\!-\!i)(k\!-\!i\!-\!1)\\
={}&\frac12(n-k)^2+\frac32(n-k)-mn+\frac12m(m+1)k+1=:f_n(m,k).
\end{align*}
Denoting $\widehat X_n:=\{(m,k,i)\in[1,\infty[^3:i\le k, mk+k-i=n\}$ and $X_n:=\widehat X_n\cap\mathbb{N}^3$ we note that
$$b_n=\max_{(m,k,i)\in X_n}f_n(m,k)\le\max_{(m,k,i)\in \widehat X_n}f_n(m,k)
\le\max_{k\in[1,n]}\left(\max_{m\in[\frac nk-1,\frac nk]}f_n(m,k)\right).$$
For fixed $k$, $f_n(m,k)$ is a convex function of $m$, and
$$f_n\left(\frac nk,k\right)=f_n\left(\frac nk-1,k\right)
%=\frac12(n-k)^2+\frac32(n-k)-\frac12\frac{n^2}{k}+\frac12n+1
=\frac{n^2}2 -nk+\frac{k^2}2+2n-\frac {3k}2-\frac{n^2}{2k}+1
:=g_n(k)\, ,$$
therefore $b_n\le\max_{k\in[1,n]}g_n(k)$. We collect some facts about $g_n$, assuming $n\ge5$. From $g'_n(k)=k-n-\frac32+\frac{n^2}{2k^2}$,
$g_n''(k)=1-\frac{n^2}{k^3}$ and $g_n'''(k)=\frac{3n^2}{k^4}$ we deduce that $g_n$ is strictly concave on $[1,n^{\frac23}]$ and strictly convex on $[n^{\frac23},n]$, and that $g_n'$ is strictly convex on $[1,n]$ with $g_n'(1)=\frac{n^2-1}2-n>0$ and
\cogg{$g_n'(n)= -1<0$} so that there is only one zero of $g_n'$ in $[1,n]$ which is the only maximizer of $g_n$ on $[1,n]$, and this maximizer must lie in the interval
$A_n:=[z_n+\frac14-\frac1{z_n},z_n+\frac14]$, where $z_n=\sqrt{\frac n2}$.
Indeed,
$$g_n'\left(z_n+\frac14\right)
=\frac12\frac{n^2}{z_n^2}\left(1+\frac1{4z_n}\right)^{-2}-n+z_n-\frac54\in\left[-\frac98,-\frac78\right]\, ,$$
where for the latter inclusion we used $(1-y)^2\le(1+y)^{-2}\le1-2y+3y^2$ for $y\in[0,1]$; we need $y=\frac1{4z_n}$. \corr{So we have shown $g_n'(\sup A_n) <0$.}
Furthermore, for $k\in A_n$ we have
$$g_n''(k)\le 1-n^2\left(z_n+\frac14\right)^{-3}=1-4z_n\left(1+\frac1{4z_n}\right)^{-3}\le1-4z_n\left(1-\frac3{4z_n}\right)=4-4z_n\, ,$$
and by the mean value theorem for some $k\in A_n$
$$g_n'\left(z_n+\frac14-\frac1{z_n}\right)=g_n'\left(z_n+\frac14\right)-\frac1{z_n}g_n''(k)\ge-\frac98+4-\frac4{z_n}\ge\frac13\, .$$
\corr{We conclude $g_n'(\sup A_n)<0<g_n'(\inf A_n )$ so that $A_n$ must contain the minimizer of $g_n$.} Now $z_n+\frac14<n^\frac23$ for $n\ge5$, and by concavity of $g_n$ on $[1,n^\frac23]$ we get
\begin{align*}\max_{k\in[1,n]}{}&g_n(k)=\max_{k\in A_n}g_n(k)\le g_n\left(z_n+\frac14\right)+\frac9{8z_n}\\
={}&\frac{n^2}2 -n\left(z_n+\frac14\right)+\frac{(z_n\!+\!\frac14)^2}2+2n-\frac {3(z_n\!+\!\frac14)}2-\frac{n^2}{2(z_n\!+\!\frac14)}+1+\frac9{8z_n}\\
\le{}&\frac{n^2}2 -2nz_n+\frac94n-\frac54z_n+\frac32\, ,
\end{align*}
where we used
$\frac{n^2}2\left(z_n+\frac14\right)^{-1}=nz_n\left(1+\frac1{4z_n}\right)^{-1}\ge nz_n\left(1-\frac1{4z_n}\right)=nz_n-\frac n4$, and $\frac{21}{32}+\frac9{8z_n}\le\frac32$ for $n\ge5$. This proves the right\corr{most} inequality in \eqref{asugl}.

Turning now to the left inequality in \eqref{asugl}, we note that for any
$(m,k,i)\in X_n$ we have $b_n\ge f_n(m,k)$. The preceding calculations suggest to choose $$k_n\!:=\!z_n\!+\!\alpha_n\in\left[z_n\!-\!\frac14,z_n\!+\!\frac34\right]\cap\mathbb{N}\, , \quad m_n:=\frac n{z_n}+\beta_n\in{}\left]\frac n{k_n}\!-\!1,\frac n{k_n}\right]\cap\mathbb{N}$$
and \corr{$i_n:=m_n k_n+k_n-n\in[1\!:\!k_n]$}. Clearly $\alpha_n\in\left[-\frac14,\frac34\right]$, \corr{and because
%the function $\frac {n\alpha}{z_n(z_n+\alpha)}$ increases with $\alpha$ on this interval,
\cogg{we have $\left\vert\frac n{k_n}-\frac n{z_n}\right\vert=2|\alpha_n|\frac{z_n}{z_n+\alpha_n}\le\frac32$
for $(\alpha_n, z_n) \in \left[-\frac14,\frac34\right]\times [1,\infty[\,$,}
we get} $\beta_n\in\left[-\frac52,\frac32\right]$. % because of $\left|\frac n{k_n}-\frac n{z_n}\right|=2|\alpha_n|\frac{z_n}{k_n}\le\frac32$.
We obtain
\begin{align*}
f_n(m_n,k_n)={}&\frac{n^2}2-nk_n-nm_n+\frac12m_n^2k_n+\frac12k_n^2+\frac{3n}2+\frac12m_nk_n
-\frac32k_n+1\\
={}&\frac{n^2}2-2nz_n+\frac94n +\gamma_nz_n+\delta_n\\
\ge{}&\frac{n^2}2-2nz_n+\frac94n-2z_n+\cogg{\frac1{16}},
\end{align*}
where we used
$$\gamma_n=
%2\alpha_n+2\alpha_n\beta_n+\frac{\beta_n^2}2+\frac{\beta_n}2-\frac32=
\frac12(\beta_n+2\alpha_n)(\beta_n+2\alpha_n+1)+\alpha_n(1-2\alpha_n)-\frac32\ge-\frac18-\frac38-\frac32=-2
$$
and, \corr{discussing behaviour for $(\alpha_n, \beta_n) \in \left[-\frac14,\frac34\right]\times \left[-\frac52,\frac32\right]$,}
%by considering the cases $\alpha_n\ge0$ (using $\beta_n(\beta_n+1)\ge-\frac14$) and $\alpha_n<0$ (using $\beta_n(\beta_n+1)\le\frac{15}4$ separately,
\corr{$$\delta_n=\frac12\alpha_n\beta_n(\beta_n+1)+\frac12\alpha_n^2- \frac 32\alpha_{\cogg{n}} +1\ge \frac 1{16}\, .$$} %\min\left(1-\frac1{128},1-\frac{7}{16}\right)\ge\frac12.$$
The proof is now complete.\ep

\begin{rem}\label{kstar} For later reference we add that $\max_{(m,k,i)\in X_n}f_n(m,k)$ is attained only in points $(m^*,k^*,i^*)$ satisfying
$$k^*\le z_n+\frac32\, .$$
To see this, let $k \,>\, z_n+\frac32$. Then, \corr{by straightforward but tedious calculations, we get} $f_n(m,k)\le g_n(k) \,<\, g_n(z_n+\tfrac32)\le \frac{n^2}2-2nz_n+\frac94n-2z_n+\corr{\frac1{16}}\le b_n$,
therefore $f_n(m,k)$ can not be maximal. \end{rem}

\begin{cly}\label{clyMain}
The DJL-conjecture is false for $n\ge7$.
Asymptotically, $p_n$ is much closer to the upper bound $s_n=\binom{n+1}{2}-4$ than to the DJL lower bound $d_n=\big\lfloor\frac{n^2}4\big\rfloor$:
\beq{boundslim}
\corr{p_n = \frac {n^2}2 + {\mathcal O}\big(n^{3/2}\big) \quad\mbox{and thus}\quad}\lim_{n\to\infty}\frac{s_n-p_n}{p_n-d_n}=0\, .\eeq
\end{cly}
\bep
For $n\in[7\!:\!11]$ counterexamples were given in \cite{BSU}, and for $n=12$ we gave a counterexample in Example~\ref{121314}.
Furthermore, we derive from \eqref{asugl}
\beq{satz32abs}
\corr{\frac {n^2}2 + {\mathcal O}\big(n^{3/2}\big) =s_n} \ge p_n\ge b_n\ge
\corr{\frac{n^2}2-(n+1)\sqrt{2n}+\frac {9}4\cogg{n} + \frac 1{16}} > d_n\, ,\eeq
where the latter inequality holds for $n\ge 13$ \corr{(again checked straightforwardly)}, showing the existence of counterexamples also for $n\ge13$. \corr{Now~\eqref{boundslim} follows immediately.} \ep

\section{Improvement of lower bounds}\label{improbo}

\subsection{{Semigroups of characteristic triples}}

Up to now, we have used in our construction a very simple matrix sequence ${\mathcal I}:=(\Ib_n)_{n\in \N}$.
This was sufficient to disprove the DJL conjecture for large $n$ and \ru_cor_1{establish} the asymptotics in \eqref{boundslim}.
Note that $b_n$ is a lower bound for the cp-rank of matrices from a subset of $\cm _{n}$,
namely for completely positive $n\times n$-matrices that have a representation as $\Ub\T\Db\Ub$,
where $\Db$ is a nonnegative diagonal matrix and $\Ub\in\mathcal{Z}$ is a binary matrix.
No longer insisting on matrices in that subset, we will be able to further increase our lower bounds for $p_n$.
So our strategy is to replace ${\mathcal I}$ by another sequence ${\mathcal J}=(\Jb_n)_{n\in\N}$ of not necessarily binary matrices,
where %for simplicity
we assume that $\Jb_n$ \ws_cor_1{has $n$ columns}, \ru_cor_1{that} all $\Jb_n$ have full column rank,
and that we know the exact values of
$\rho^{{\mathcal J}}_n:=\ran2\Jb_n$, not just lower bounds, with $\rho^{\mathcal J}_n>n$ for at least one $n$.
Then we let $S^{\mathcal{J}}$ be the semigroup generated by the set of characteristic triples
$\{(n,n,\rho^{\mathcal J}_n):n\in\mathbb{N}\}$, and define
\beq{rhoundbJ}b_n^{\mathcal{J}}:=\max\cogg{\big\{}\corrr{c_3 : c_1}=n\, ,\; c\in S^{\mathcal{J}}\cogg{\big\}}.\eeq
We recall that $S^{\mathcal I} =S$ and $b_n^{\mathcal I}=b_n$ from~\reff{rhoundb} in this notation,
and of course $\rho^{\mathcal I}_n=n$. Further, for all such  ${\mathcal J}$, from considering $\corrr{(c\rad c')_2}$,
we deduce that any $c\in S^{\mathcal J}$ is $\rad$-\cogg{irreducible} if and only if $\corrr{c_1=c_2}$.
In other words,  $(n,n,\rho)\in S^{\mathcal J}$ if and only if $\rho = \rho_n^{\mathcal J}$.
\cogg{Clearly we may again infer that there is $\Mb=\Ub\T\Db\Ub\in\cm _n$ satisfying $\cpr \Mb\ge b_n^{\mathcal J}$,
where $\Db$ is a nonnegative diagonal matrix and $\Ub$ is an element of the subsemigroup of $(\mathcal{Z},\rad)$
generated by $\mathcal J$. Such $\Ub$ can be found \corr{as follows: take} some maximizing characteristic triple $c\in S^{\mathcal{J}}$
satisfying $\corrr{c_1}=n$ \corr{and} $\corrr{c_3}=b_n^{\mathcal{J}}$ (there may be more than one maximizing characteristic triple);
use some factorization of $c$ as a product of generators (again there may be more than one such factorization),
say $c=\corrr{c^{(1)}\rad\cdots\rad c^{(k)}}$, for some $k\in\mathbb{N}$; and define $\Ub:=\Jb_{\corrr{c^{(1)}_1}}\rad\cdots\rad\Jb_{\corrr{c^{(k)}_1}}$.}

\ws_cor_1{Note that we could have included several building blocks for a given $n$, but did not for the following reason:
Adding $(n,k,\rho)$ to the set of generators only pays if $\rho>\rho'$ for all $(n,k,\rho')\in S^{\mathcal J}$, \corrr{by~\reff{monoprop}.} So we included exactly one generator for each $n$, if $k=n$, and no other generators, since in the case $k<n$ we were not able to directly come up with generators that would beat those already in $S^{\mathcal J}$.}


The next \corr{result} is about the increase $b_n^{\mathcal{J}}-b_n$ of the lower bound that we may expect in the case that we have certain bounds for $\rho^{\mathcal J}_n$.
\begin{lem}\label{alphabeta}
Assume that for some $\alpha,\beta>0$ we have $(\alpha+1) n-\beta\le\rho^{\mathcal J}_n\le(\alpha+1) n$ for $n\in\mathbb{N}$. Then $b_n^{\mathcal{J}}$ satisfies
$$\alpha n-\beta\left(\sqrt{\frac n2}+\frac32\right)\le b_n^{\mathcal{J}}-b_n\le \alpha n,$$
with $b_n$ as defined in \eqref{rhoundb}.
\end{lem}
\bep
We are going to show by induction on $k$ that\\
a) for each $c=(n,n+1-k,\rho)\in S^{\mathcal{J}}$ there is $c'=(n,n+1-k,\rho')\in S$ satisfying
$$\rho\le \rho'+\alpha n,$$
b) for each $c'=(n,n+1-k,\rho')\in S$ there is $c=(n,n+1-k,\rho)\in S^{\mathcal{J}}$ satisfying
$$\rho'+\alpha n-\beta k\le \rho.$$
Both assertions are true for $k=1$, with $\rho=\rho^{\mathcal J}_n$ and $\rho'=n$.
Next we assume a) proved up to $k$, and we use that
$c=(n,n-k,\rho)\in S^{\mathcal{J}}$ has for some $i\in[1:n-1]$ a representation $c=(i,i,\rho^{\mathcal J}_i)\rad\bar c$, where $\bar c=(n-i,n-i-k+1,\bar\rho)\in S^{\mathcal{J}}$.
By assumption, there is $\bar c'=(n-i,n-i-k+1,\bar\rho')\in S$ such that
$$\bar\rho\le \bar\rho'+\alpha(n-i),$$
and then $c':=(i,i,i)\rad\bar c'=(n,n-k,\rho')\in S$ satisfies
$$\rho-\rho'=\rho^{\mathcal J}_i-i+\bar\rho-\bar\rho'\le \alpha i+\alpha(n-i)=\alpha n.$$
Next, assuming b) proved up to $k$, we use that
$c'=(n,n-k,\rho')\in S$ has for some $i\in[1:n-1]$ a representation $c'=(i,i,i)\rad\bar c'$, where $\bar c'=(n-i,n-i-k+1,\bar\rho')\in S$.
By assumption, there is $\bar c=(n-i,n-i-k+1,\bar\rho)\in S^{\mathcal{J}}$ such that
$$\bar\rho'+\alpha (n-i)-\beta\cogg{k}\le \bar\rho,$$
and then $c':=(i,i,\rho^{\mathcal J}_i)\rad\bar c=(n,n-k,\rho)\in S^{\mathcal{J}}$ satisfies
$$\rho-\rho'=\rho^{\mathcal J}_i-i+\bar\rho-\bar\rho'\ge \alpha i-\beta+\alpha(n-i)-\beta \cogg{k}=\alpha n-\beta \cogg{(k+1)}.$$
Now we use a) to obtain % and b) to obtain
$$b_n^{\mathcal{J}}=\max\limits_{r,\rho}\lk \rho:(n,r,\rho)\in S^{\mathcal{J}}\rk\le
\max\limits_{r,\rho'}\lk\rho'+\alpha n:(n,r,\rho')\in S\rk=b_n+\alpha n$$
and, using b) and Remark~\ref{kstar},
\begin{align*}
b_n+\alpha n-\beta\left(\sqrt{\frac n2}+\frac32\right)
\le{}
&\max\limits_{\cogg{k},\rho'}\lk \rho'+\alpha n-\beta k:(n,n+1-k,\rho')\in S\rk\\
\le{}
&
\max\limits_{\cogg{k},\rho}\lk \rho:(n,n+1-k,\rho)\in S^{\mathcal{J}}\rk = b_n^{\mathcal{J}}\, .
\end{align*}
Hence the results.\ep

\corr{So by this method we always obtain an improvement which increases linearly in $n$, but we cannot hope for much more. The next theorem makes this more precise, and also provides a construction principle for such an improving sequence $\mathcal J$:}

\begin{thm}\label{uff}
\corr{Suppose} %We meet the hypothesis of Lemma~\ref{alphabeta} whenever
we choose ${\mathcal J}=(\Jb_n)_{n\in\N}$ as follows:
\begin{itemize}
  \item Fix $n_0\in \N$ and select $\Jb_n\in\mathcal{Z}$ with full column rank (and $\rho^{\mathcal J}_n=\ran2\Jb_n$) for \cogg{all} $n\in[1\!:\!n_0]$\cogg{, with $\rho^{\mathcal J}_k>k$ for at least one $k\in[1\!:\!n_0]$}.
  \item Let $k_0:=\min\cogg{\Big\{}n\in[1\!:\!n_0]:\frac{\rho^{\mathcal J}_n}n\ge\frac{\rho^{\mathcal J}_\ell}\ell\textup{ for all }\ell\in[1\!:\!n_0]\cogg{\Big\}}$.
  \item Write any $n>n_0$ as $n=ak_0+b$, where $n_0-k_0< b\le n_0$. Abbreviating
$$q\odot \Ab := \Ab\oplus \Ab \oplus \dots \oplus \Ab$$
for $q$ such $\oplus$-operands $\Ab$, %on the right-hand side,
define $\Jb_n=\left (a\odot\Jb_{k_0}\right ) \oplus\Jb_b\in\mathcal{Z}$, which is a matrix of full column rank by Lemma~\ref{ran2rec},
  and let $\rho^{\mathcal J}_n=\ran2\Jb_n=a\rho^{\mathcal J}_{k_0}+\rho^{\mathcal J}_b$.
 \end{itemize}
\corr{Then $b_n^{\mathcal{J}}-b_n = \alpha n + {\mathcal O}(\sqrt n)$ for some $\alpha >0$ and thus
\beq{uffabs} \frac {n^2}2 + \frac n2 -4 \ge p_n \ge \frac {n^2}2 - \sqrt 2 n^{3/2} + \gamma n +  {\mathcal O}(\sqrt n)\eeq
for some $\gamma > \frac 94$ depending on the first $n_0$ matrices $\Jb_n$, $n\in [1\!:\!n_0]$.}
\end{thm}

\bep
With $\alpha':=\frac{\rho^{\mathcal J}_{k_0}}{k_0}\cogg{>1}$ we have $\rho^{\mathcal J}_n\le\alpha' n$ for $n\in[1\!:\!n_0]$ by the definition of $\alpha'$, and
$\rho^{\mathcal J}_n=a\rho^{\mathcal J}_{k_0}+\rho^{\mathcal J}_b\le a\alpha' k_0+\alpha' b=\alpha'(ak_0+b)=\alpha' n$ also for $n>n_0$.
With $\beta:=\max\cogg{\big\{}\alpha' n-\rho^{\mathcal J}_n:n\in[1\!:\!n_0]\cogg{\big\}\ge\alpha'-1>0}$ we have $\rho^{\mathcal J}_n\ge\alpha' n-\beta$ for $n\in[1\!:\!n_0]$, and
$\rho^{\mathcal J}_n=a\rho^{\mathcal J}_{k_0}+\rho^{\mathcal J}_b\ge a\alpha' k_0+\alpha' b-\beta=\alpha' n-\beta$ also for $n>n_0$. So the hypothesis of Lemma~\ref{alphabeta} is fulfilled with
$\alpha:=\alpha'-1$ and $\beta$, \corr{and the results follow.} \ep

\subsection{{New building blocks and better bounds}}

The following example shows the construction of a particular sequence $\mathcal{J}=(\Jb_n)_{n\in\N}$, and reports on the lower bounds $b_n^{\mathcal{J}}$ obtained from this sequence.

\renewcommand\arraystretch{1}
\begin{table}[ht]
\caption{Matrices $\Jb_n$ from Example~\ref{ex2}. See text.}
\label{tab:goodJn}
%\begin{minipage}{\textwidth}
%\renewcommand{\footnoterule}{}
%\renewcommand{\thefootnote}{\alph{footnote}}
\begin{center}
\begin{tabular}{| c| c|c|| c| c|c|| c| c|c|}
\hline
   $n$     & $\Jb_n$    & $\rho^{\mathcal J}_n$   & $n$ & $\Jb_n$  & $\rho^{\mathcal J}_n$ & $n$ & $\Jb_n$  & $\rho^{\mathcal J}_n$\\[0.2em]
\hline %\hlin
\multirow{2}{15mm}{$k\in[1\!:\!5]$} & \multirow{2}{5mm}{~$\Ib_k$} & \multirow{2}{5mm}{~$k$} & $9$ & $\Jb_9$  & $26$ & $14$ & $\Jb_{14}$                & $80$ \\
\cline{4-9}
                                &                             &                         &$10$ & $\Jb_9\oplus \Jb_1$      & $27$ & $15$ & $\Jb_{15}$                & $95$ \\
\hline
$6$             & $\Jb_6$                     & $8$                       &$11$ & $\Jb_{11}$          & $32$ & $26$ & $\Jb_{14}\oplus \Jb_{12}$ & $130$\\
\hline
$7$             & $\Jb_7$                     & $14$                      &$12$ & $\Jb_{12}$          & $50$ &\multirow{2}{15mm}{\small ~~$k+15$, $k\!\in\!\mathbb{N}\!\setminus\! \{11\}$} & \multirow{2}{13mm}{$\Jb_{k}\oplus \Jb_{15}$} & \multirow{2}{12mm}{$\rho^{\mathcal J}_k+95$}\\
\cline{1-6}
$8$             & $\Jb_8$                     & $18$                      &$13$ & $\Jb_{13}$          & $65$ &    &                           &     \\
\hline\end{tabular}
\end{center}
%\normalsize \footnotetext[1]{~see Step~1} \footnotetext[2]{~see Step~2} \footnotetext[3]{~see Step~3}
%\end{minipage}
\end{table}

\renewcommand\tabcolsep{3pt}
\renewcommand\arraystretch{1.2}
\begin{table}[ht]
\caption{Several bounds for $p_n$, where $b_n^{\mathcal{J}}$ is a lower bound for $\ran2 \Ub_n $; see text.}
\label{tab:goodUn}
\begin{center}
\scriptsize{\begin{tabular}{| c| c| c| c|c|c|| c| c|c| c| c|c|}
\hline
   $n$     & $d_n$  &   $b_n$     & $b_n^{\mathcal{J}}$  & $\Ub_n$ & $s_n$ &
   $n$     & $d_n$  &   $b_n$     & $b_n^{\mathcal{J}}$  & $\Ub_n$ & $s_n$\\[0.2em]
\hline %\hline
$6$  & 9  &  9 &  9 & $\Jb_{3}\rad \Jb_{3}$   & 17        & $40$  & 400&526 & 664 & $\Jb_{13}^{\ssrad 2}\rad \Jb_{14} $   & 816                   \\
$7$  & 12 & 12 &$\bf{14}$& $\Jb_{7}$ & 24                 & $45$  & 506&681 &871 & $\Jb_{15}^{\ssrad 3}$ & 1031                               \\
$8$  & 16 & 16 & 18 & $\Jb_{8}$ & 32                      & $50$  & 625&856 & 1043  & $\Jb_{5}\rad \Jb_{15}^{\ssrad 3}$ & 1271                                    \\
$9$  & 20 & 20 & 26 & $\Jb_{9}$ & 41                      & $55$  & 756& 1051 & 1277  & $\Jb_{13}\rad \Jb_{14}^{\ssrad 3}$ & 1536                                    \\
$10$ & 25 & 25 & 28 & $\Jb_{3}\rad \Jb_{7}$ & 51          & $60$ & 900& 1270 & 1553  & $\Jb_{15}^{\ssrad 4}$ & 1826                     \\
$11$ & 30 & 30 & 35 &$\Jb_{2}\rad \Jb_{9}$     & 62       & $65$ & 1056& 1510 & 1781  &$\Jb_{5}\rad \Jb_{15}^{\ssrad 4}$  & 2141                   \\
$12$ & 36 &$\bf{37}$& 50 &$\Jb_{12}$     & 74             & $70$ & 1225& 1771 & 2086 &$\Jb_{14}^{\ssrad 5}$&  2481                          \\
$13$ & 42 & 44 & 65 &$\Jb_{13}$     & 87                  & $80$ & 1600& 2357 & 2726 &$\Jb_{16}^{\ssrad 5}$&  3236                               \\
$14$ & 49 & 52 & 80 &$\Jb_{14}$     & 101                 & $90$ & 2025& 3036 & 3505 &$\Jb_{15}^{\ssrad 6}$&    4091                             \\
$15$ & 56 & 61 & 95 &$\Jb_{15}$    & 116                  & $100$ & 2500& 3800 & 4290  &$\Jb_{14}^{\ssrad 5}\rad \Jb_{15}^{\ssrad 2}$ & 5046                           \\
$16$ & 64 & 70 & 96 & $\Jb_{16}$    & 132                 & $120$ & 3600& 5601 & 6241 & $\Jb_{15}^{\ssrad 8}$    & 7256                               \\
$17$ & 72 & 80 & 110 & $\Jb_{2}\rad \Jb_{15}$    & 149    & $140$ & 4900& 7758 & 8478  & $\Jb_{15}^{\ssrad 4}\rad \Jb_{16}^{\ssrad 5}$& 9866                \\
$18$ & 81 & 91 & 125 & $\Jb_{3}\rad \Jb_{15}$    & 167    & $160$ & 6400&  10285 & 11076  & $\Jb_{16}^{\ssrad 10}$    & 12876                \\
$19$ & 90 & 102& 140& $\Jb_{4}\rad \Jb_{15}$    & 186     & $180$ & 8100&  13176& 14065 & $\Jb_{15}^{\ssrad 12}$    & 16286                 \\
$20$ & 100& 114& 155& $\Jb_{5}\rad \Jb_{15}$& 206         & $200$ & 10000&16436& 17366& $\Jb_{15}^{\ssrad 8}\rad \Jb_{16}^{\ssrad 5}$& 20096                     \\
$21$ & 110& 127& 172& $\Jb_{6}\rad \Jb_{15}$  & 227       & $250$ & 15625&26203& 27261& $\Jb_{18}\rad \Jb_{22}\rad \Jb_{30}^{\ssrad 7}$  & 31371                   \\
$22$ & 121& 140& 192& $\Jb_{7}\rad \Jb_{15}$ & 249        & $300$ & 22500&38305& 39736& $\Jb_{30}^{\ssrad 10}$ & 45146                    \\
$23$ & 132& 155& 210& $\Jb_{8}\rad \Jb_{15}$ & 272        & $350$ & 30625&52754& 54495& $\Jb_{20}\rad \Jb_{30}^{\ssrad 11}$ & 61421                    \\
$24$ & 144& 171& 232& $\Jb_{9}\rad \Jb_{15}$ & 296        & $400$ & 40000&69562& 71591& $\Jb_{30}^{\ssrad 3}\rad \Jb_{31}^{\ssrad 10}$& 80196 \\
$25$ & 156& 187& 247& $\Jb_{10}\rad \Jb_{15}$ & 321       & $450$ & 50625&88741& 91141& $\Jb_{30}^{\ssrad 15}$ & 101471                   \\
$26$ & 169& 204 & 273 & $\Jb_{13}\rad \Jb_{13}$ & 347     & $500$ & 62500& 110291 & 112860 & $\Jb_{20}\rad \Jb_{30}^{\ssrad 16}$    & 125246               \\
$27$ & 182& 222 & 300 & $\Jb_{13}\rad \Jb_{14}$ & 374     & $550$ & 75625& 134221 & 137061 & $\Jb_{30}^{\ssrad 8}\rad \Jb_{31}^{\ssrad 10}$& 151521              \\
$28$ & 196& 241 & 328 & $\Jb_{14}\rad \Jb_{14}$ & 402     & $600$ & 90000& 160534 & 163571 & $\Jb_{30}^{\ssrad 20}$    & 180296              \\
$29$ & 210& 260& 356& $\Jb_{14}\rad \Jb_{15}$  & 431      & $650$ & 105625& 189249& 192390& $\Jb_{30}\rad \Jb_{31}^{\ssrad 20}$    & 211571                 \\
$30$ & 225& 280& 385& $\Jb_{15}\rad \Jb_{15}$& 461        & $700$ & 122500& 220357& 223592& $\Jb_{32}^{\ssrad 17}\rad \Jb_{33}^{\ssrad 2}\rad \Jb_{45}^{\ssrad 2}$& 245346                    \\
$31$ & 240& 301& 400& $\Jb_{15}\rad \Jb_{16}$  & 492      & $734$ & 134689& 242873& 246353& $\Jb_{32}\rad\Jb_{33}\rad\Jb_{39}\rad \Jb_{45}^{\ssrad 14}$  & 269741                   \\
$32$ & 256& 323& 416& $\Jb_{16}^{\ssrad 2}$ & 524       & $800$ & 160000& 289771& 293751& $\Jb_{35}\rad \Jb_{45}^{\ssrad 17}$ & 320396                    \\
$33$ & 272& 345& 443& $\Jb_{3}\rad \Jb_{15}^{\ssrad 2}$& 557&$850$ & 180625& 328085& 332428& $\Jb_{44}^{\ssrad 5}\rad \Jb_{45}^{\ssrad 14}$ & 361671  \\
$34$ & 289& 368& 472& $\Jb_{4}\rad \Jb_{15}^{\ssrad 2}$& 591&$900$ & 202500& 368803& 373521& $\Jb_{45}^{\ssrad 20}$ & 405446  \\
$35$ & 306& 392& 501& $\Jb_{5}\rad \Jb_{15}^{\ssrad 2}$& 626&$1000$& 250000& 457489& 462760& $\Jb_{45}^{\ssrad 12}\rad \Jb_{46}^{\ssrad 10}$ & 500496  \\
\hline
\end{tabular}}
\end{center}
\end{table}

\begin{exa}\label{ex2}
We let $n_0=26$. The construction of $\Jb_n$ for $n\in[1\!:\!n_0]$ will be divided into 3 steps, and $\Jb_n$ for $n>n_0$ will be constructed in a fourth step.



\textbf{Step 1:} We start with elementary {\em building blocks} $\Jb_n:=\Ib_n$ for $n\in[1\!:\!5]$, and we add four more building blocks $\Jb_n\in\mathcal{Z}$ satisfying $\ran2 \Jb_n>n$. These we get by looking for copositive matrices having many zeroes, in particular we employ, using the notation $\Cb(\a)$ with $\a\in\R^n$ from~\cite{BSU},
%~\vspace*{-.3cm}{\scriptsize\begin{align*}
~\vspace*{-.15cm}{\small\begin{align*}
\Sb_7={}&\Cb([-153,127,-27,-27,127,-153,162]\T)\,,\\
\Sb_9={}&\Cb([-1056, 959, -484, 231, 231, -484, 959, -1056, 1089]\T)\, ,\\
\Sb_{11}={}&\Cb([32, 18, 4, -24, -31, -31, -24, 4, 18, 32, 32]\T)\, ,\quad\mbox{and}\\
\Sb_{15}={}&\Cb([a,a,b,a,c,b,a,a,b,c,a,b,a,a,d]\T)\, ,\ \textup{where}\ [a,b,c,d]=[-2609, 1803, 4009, 7318]\, .
%\Sb_{15}={}&\Cb([-2609, -2609, 1803, -2609, 4009, 1803, -2609, -2609, 1803, 4009, -2609, 1803, -2609, -2609, 7318]\T)\, .
\end{align*}}%
\vspace*{-.75cm}\\
Let the rows of $\Jb_7\in\mathbb{R}^{14\times 7}$ (resp.~$\Jb_9\in\mathbb{R}^{27\times 9}$,
$\Jb_{11}\in\mathbb{R}^{33\times 11}$, $\Jb_{15}\in\mathbb{R}^{360\times 15}$) be the zeroes of $\Sb_7$, (resp.~$\Sb_9$, $\Sb_{11}$, $\Sb_{15})$.
Those matrices all have full column rank and satisfy $\ran2 \Jb_7 =14$, $\ran2 \Jb_9 =26$, $\ran2 \Jb_{11} =32$
and $\ran2 \Jb_{15} =95$.

\textbf{Step 2:} Now we delete some rows to close the gaps in column numbers, i.e., consider $n\in \lk 6,8,12,13,14\rk$. Generally speaking, if
$\Ub\in\mathbb{R}^{k\times n}$ collects in its rows all the zeroes of $\Sb\in\mathbb{R}^{n\times n}$, we define for a subset $N\subseteq[1\!:\!n]$ the complement $N'=[1\!:\!n]\setminus N$ and put $K:=\{\ell\in[1\!:\!k]:u_{\ell,i}=0\textup{ for all }i\in N\}$. Finally, we abbreviate by $\Phi_N(\Ub):=\Ub_{K\times N'}$, so that the rows of $\Phi_N(\Ub)$ are the zeroes of the matrix $\Sb_{N'\times N'}$. The motivation is that if $\Ub$ has full rank and large two-rank, then in lucky cases the same will be true for $\Phi_N(\Ub)$ for small sets $N$. Indeed, $\Jb_6:=\Phi_{\{1\}}(\Jb_7)\in\mathbb{R}^{8\times 6}$, $\Jb_8:=\Phi_{\{1\}}(\Jb_9)\in\mathbb{R}^{18\times 8}$, $\Jb_{12}:=\Phi_{\{1,2,3\}}(\Jb_{15})\in\mathbb{R}^{60\times 12}$, $\Jb_{13}:=\Phi_{\{1,2\}}(\Jb_{15})\in\mathbb{R}^{108\times 13}$ and $\Jb_{14}:=\Phi_{\{1\}}(\Jb_{15})\in\mathbb{R}^{192\times 14}$
have all full column rank and satisfy $\ran2 \Jb_6 =8$, $\ran2 \Jb_8 =18$, $\ran2 \Jb_{12} =50$, $\ran2 \Jb_{13} =65$ and $\ran2 \Jb_{14} =80$.

\textbf{Step 3:} %\cor{The case $n=10$ is still not treated. Therefore,} using Lemma~\ref{ran2rec}(c),
We further define $\Jb_{10}:=\Jb_9\oplus \Jb_1\in\mathbb{R}^{28\times 10}$, satisfying $\rank \Jb_{10} =10$ and $\ran2 \Jb_{10} =27$.
For $n\in[16\!:\!25]$ we define $\Jb_n:=\Jb_{n-15}\oplus \Jb_{15}$, and, deviating from the latter pattern, we finally let $\Jb_{26}:=\Jb_{12}\oplus \Jb_{14}$,
because $\ran2\Jb_{12}+\ran2\Jb_{14}=130>127=\ran2\Jb_{11}+\ran2\Jb_{15}$, and this completes the construction of $\Jb_n$ and $\rho^{\mathcal J}_n:=\ran2\Jb_n$ for $n\in[1\!:\!26]$. We remark that the matrices $\Jb_i,i\in[7\!:\!11]$ have also been used in the paper~\cite{BSU} to provide the first counterexamples to the DJL conjecture.

\textbf{Step 4:} We compute $k_0:=\min\lk n\in[1\!:\!26]:\frac{\rho^{\mathcal J}_n}n\ge\frac{\rho^{\mathcal J}_\ell}\ell\textup{ for all }\ell\in[1\!:\!26]\rk=15$.
Any $n>26$ is now written as $n=15a+b$, where $11< b\le 26$, and accordingly we define $\Jb_n=a\odot\Jb_{15}\oplus\Jb_b\in\mathcal{Z}$,
and let $\rho^{\mathcal J}_n=\ran2\Jb_n=95a+\rho^{\mathcal J}_b$. This completes the construction of the full sequence $\mathcal{J}$, a summary of which can be found in Table~\ref{tab:goodJn}.

We note without proof that the sequence $(\rho^{\mathcal J}_n)_{n\in\mathbb{N}}$ satisfies
$\rho^{\mathcal J}_n\ge\rho^{\mathcal J}_i+\rho^{\mathcal J}_{n-i}$ for any $i,n\in\mathbb{N}$ such that $i<n$,
and so our construction can be seen as picking from the semigroup \corrr{generated by
$\{\Jb_n:n\in[1\!:\!26]\}$ with binary operation $\oplus$, for each column dimension, one of the matrices of highest \cogg{two-rank}, in view of Lemma~\ref{ran2rec}(c)}.
The values of $b_n^{\mathcal{J}}$ for some $n\in[1\!:\!1000]$, together with other bounds and matrices $\Ub_n$ achieving
$\ran2 \Ub_n\ge b_n^{\mathcal{J}}$ are given in Table~\ref{tab:goodUn}. The data for all other $n\in [36\!:\!999]$
are available from the authors upon request. See also Figure~\ref{fig:bounds}.
\cogg{The matrices $\Ub_n$ listed in Table~\ref{tab:goodUn} have been obtained as outlined in the beginning of this section, and are therefore
in general not unique. For instance, we could also have chosen $\Ub_{11}\!=\!\Jb_4\rad\Jb_7$,
because $(2,2,2)\rad(9,9,26)\!=\!(4,4,4)\rad(7,7,14)\!=\!(11,10,35)$. Note that $b_{10}^{\mathcal{J}}\!=\!28$ and $b_{11}^{\mathcal{J}}\!=\!35$
provide better lower bounds for $p_{10}$ and $p_{11}$ than $27$ and $32$, the ones given in \cite{BSU}.
As may be seen from the right half of Table~\ref{tab:goodUn}, the structure
of maximizers of \eqref{rhoundbJ} is more complicated than the structure of maximizers of \eqref{rhoundb}. Indeed, there is no simple analogue of
Lemma~\ref{twofactors}, since a maximizing characteristic triple from $S^{\mathcal{J}}$ may need more than 2 different generators in any of its
factorizations. In the range $[1\!:\!1000]$ we found that at most 4 different generators always suffice, and that 4 are necessary in 4 cases, the
smallest of them~being~$n=734$.}

In order to get a grip on the asymptotic behavior of $(b_n^{\mathcal{J}})_{n\in\mathbb{N}}$ we compute
$\alpha:=\frac{\rho^{\mathcal J}_{15}}{15}-1=\frac{16}3$ and
$\beta:=\max\cogg{\big\{}(\alpha+1) n-\rho^{\mathcal J}_n:n\in[1\!:\!26]\cogg{\big\}}=11(\alpha+1)-\rho^{\mathcal J}_{11}=\frac{113}3$. Then $(\alpha+1) n-\beta\le\rho^{\mathcal J}_n\le(\alpha+1) n$ holds for $n\in\mathbb{N}$,
and combining Theorem~\ref{ninequarters}, Lemma~\ref{alphabeta} \corr{and Theorem~\ref{uff}}, we get
\cogg{$$-\frac{119}6\sqrt{2n}-\frac{903}{16}
\le b_n^{\mathcal{J}}-\frac{n^2}2+ n\sqrt{2n}-\frac{91}{12}n\le -\frac58\sqrt{2n}+\frac32\, .$$}
%$$-\tfrac{119}3\sqrt{\tfrac n2}-\cogg{\tfrac{903}{16}}
%\le b_n^{\mathcal{J}}-\tfrac12{n^2}+ \sqrt{2}\, n^{\frac32}\cogg{-\tfrac{91}{12}n}\le \cogg{-}\tfrac54\sqrt{\tfrac n2}+\tfrac32\, .$$

%\ws_cor_1{Das folgende besser in den response letter verschieben: Finally, note that a further improvement could be made. It is well known that non-singular $5\times5$-matrices with cp-rank %$6$ exist~\cite[Thm.~3.4]{Shak13b}, although we do not know a concrete instance (the ones in the literature are singular). Suppose there were a $5\times5$-matrix $\Mb=\Ub\T\Ub$ with %$\rank\,\Mb=5$
%and $\cpr \Mb=6$, that additionally satisfies $\Ub\in\mathcal{Z}$ (we do not know if that can happen). In that case
%employing a building block $\Jb_5=\Ub$ would increase the lower bounds $b_n^{\mathcal{J}}$ by one for $n\in\{20,35,50,65\}$, but have no other effect.
%We thank an anonymous referee for pointing out that improvement.}

\begin{figure}[h]
\centering
\includegraphics[scale=1]{boundsdjlcolor_gray.eps}%{boundsdjl.eps}%{bounds_djl.pdf}
\caption{Three lower bounds on $p_n$ compared to the best known upper bound $s_n$ } \label{fig:bounds}
\end{figure}
\end{exa}
\section{Conclusions}
Summarizing our findings regarding the DJL conjecture: it is true for $n = 5$~\cite{Shak13a}; it is false for $n\ge 7$~(see~\cite{BSU} for $n\le11$); and it is still unresolved for $n=6$, despite recent efforts to reduce the gap between the bounds for $p_6$~\cite{Shak13b}; {see also~\cite{Hild14a}.}

%\vskip0.3cm
%\noindent


\begin{appendix}

\section{Constructing a $12\times 12$-matrix of cp-rank 37.}

Here we {explicitly construct} a matrix $\Mb\in\cm_{12}$ with $\cpr \Mb=37$, as announced in Example~\ref{121314}.
\renewcommand\arraystretch{0.6}
Let the matrix $\Ub=\left[
           \begin{array}{c}
            \! \Ub_0\! \\
            \! \Ub_1\! \\
            \! \Ub_2\! \\
            \! \Ub_3\! \\
           \end{array}
         \right]\in\mathbb{R}^{64\times12}$
be a rearrangement of the rows of $\Ib_4^{\srad3}$ (which are binary vectors with exactly one unit entry in each of the three four-entry blocks), satisfying
$\Ub_i\in\mathbb{R}^{k_i\times12}$, where $(k_0,\ldots,k_3)=(27,27,9,1)$, and
$\Ub_i(\e_1+\e_5+\e_9)=i\,\pmb{\eta}_{k_i}$ for $i\in[0:3]$.
\renewcommand\arraycolsep{3pt}
Define the completely positive matrix\\[.3cm]
$\Mb:=6\Ub_1\T\Ub_1+6\Ub_2\T\Ub_2+\Ub_3\T\Ub_3=
\left[
           \begin{array}{cccccccccccc}
91&0&0&0&19&24&24&24&19&24&24&24\\0&42&0&0&24&6&6&6&24&6&6&6\\0&0&42&0&24&6&6&6&24&6&6&6\\0&0&0&42&24&6&6&6& 24&6&6&6\\19&24&24&24&91&0&0&0&19&24&24&24\\24&6&6&6&0&42&0&0&24&6&6&6\\24&6&6&6&0&0&42&0&24&6&6&6\\24&6&6&6& 0&0&0&42&24&6&6&6\\19&24&24&24&19&24&24&24&91&0&0&0\\24&6&6&6&24&6&6&6&0&42&0&0\\24&6&6&6&24&6&6&6&0&0&42& 0\\24&6&6&6&24&6&6&6&0&0&0&42           \end{array}
         \right],$\\[.3cm]
(observe \ru_cor_1{that} $\Mb = \Ib_3\otimes \Ab + (\Eb_3-\Ib_3)\otimes \Bb$ where $\Ab$ and $\Bb$ are the upper-left and upper-right corner $4\times 4$ blocks) and let $\mathcal{K}_m=\{(r,s)\in[1:12]^2:r<s,M_{rs}=m\}$ for $m\in\{6,19,24\}$. Clearly we have $|\mathcal{K}_6|=27$ and $|\mathcal{K}_{24}|=18$.
Furthermore consider, by analogy of the construction in Lemma~\ref{ran2rec}, the copositive matrix
$$\Sb=\left[\begin{array}{ccc}3\Eb_4-\Ib_4&-\Eb_4&-\Eb_4\\
-\Eb_4&3\Eb_4-\Ib_4 &-\Eb_4\\
-\Eb_4&-\Eb_4&3\Eb_4-\Ib_4
\end{array}\right] = \Ib_3\otimes (3\Eb_4-\Ib_4) - (\Eb_3-\Ib_3)\otimes \Eb_4 = 4\Ib_3\otimes \Eb_4 - \Ib_{12}-\Eb_{12}$$
which has exactly the rows of $\frac13\Ub$ as zeroes (indeed, $\Sb\u = \pmb{\eta}_{12}- \u$ for the rows $\u\T$ of $\Ub$). Therefore $\langle \Mb,\Sb\rangle=0$. \ws_cor_1{Next}, form the (not copositive) $4\times 4$ matrix
$\Cb= \frac 1{22}\left[\begin{array}{cccc}
5&-6&-6&-6\\-6&5&5&5\\-6&5&5&5\\-6&5&5&5
\end{array}\right]$ and
$\bar\Sb=\Sb+ \Eb_3\otimes \Cb$. %= \Ib_3\otimes (4\Eb_4 -\Ib_4) +\Eb_{3}\otimes (\Cb-\Eb_4)\, .$$
%Here $\Sb$ is one of those copositive matrices having , and
By computing all stationary points for the problem
$\min\limits_{\u\in \Delta_{12}} \u\T{\bar\Sb}\u$, it is straightforwardly checked that also $\bar \Sb$ is copositive.
We further note $\langle \Mb,\bar\Sb\rangle=0 + 3\langle \Ab,\Cb \rangle + 6\langle \Bb,\Cb\rangle =\frac{261}{22}$ and $\u\T\bar\Sb\u=\frac{45}{22}$ for any row $\u\T$ of $\Ub_0$. Now consider any \cogg{cp-factorization}
$\Mb=\Ub\T\textup{Diag}(\x)\Ub=\sum_{i=1}^{64}x_i\u_i\u_i\T$ with $\x\in\mathbb{R}_+^{64}$ and denote
$\Mb_{123}:=\Mb-\sum_{i=1}^{27}x_i\u_i\u_i\T$. As also $\Mb_{123}$ is completely positive, we have $\langle \Mb_{123},\bar \Sb\rangle \ge0$ and thus
$\sum_{i=1}^{27}x_i\le \frac{261}{45}<6$, so
we have $(\Mb_{123})_{rs}>0$ for all
$(r,s)\in\mathcal{K}_6$. Now, for any $(r,s)\in\mathcal{K}_6$, there is exactly one row $\u\T$ of
$\left[\begin{array}{c}\! \Ub_1\! \\ \! \Ub_2\! \\ \! \Ub_3\! \\ \end{array}\right]$
satisfying $u_ru_s>0$ (which must be a row of $\Ub_1$, e.g. $(\e_1+\e_6+\e_{10})\T$ for $(6,10)\in\mathcal{K}_6$). This row $\u\T$ moreover satisfies $u_\rho u_\sigma=0$ for every $(\rho,\sigma)\in\mathcal{K}_6\setminus \lk (r,s)\rk $.
As $|\mathcal{K}_6|=27$, the number of rows in $\Ub_1$,
we conclude that $0<x_i\le6$ must hold for all $i\in[28:54]$, with $x_i<6$ for some $i\in[28:54]$, if $x_i>0$ for some $i\in[1:27]$.
\ws_cor_1{Now} consider any $(r,s)\in\mathcal{K}_{24}$. First we note $(\u_{64}\u_{64}\T)_{rs}=0$ because $\mathcal{K}_{24}\cap \lk 1,5,9\rk^2 = \emptyset$.
Moreover, by a similar reasoning also
$$\left(\sum_{i=1}^{27} x_i\u_i\u_i\T\right )_{rs} =0\, .$$
Further,  there are exactly three rows $\u\T$ of $\Ub_1$ such that $u_ru_s>0$ (in case $(r,s)=(1,6)$, these are $(\e_1+\e_6+\e_{10})\T$,
$(\e_1+\e_6+\e_{11})\T$ and $(\e_1+\e_6+\e_{12})\T$) so that we arrive by above observations at
$$\left(\sum_{i=28}^{54} x_i\u_i\u_i\T \right)_{rs} \le 3\cdot 6 =18\, . $$
Next denote $\Mb_{23}:=\Mb_{123}-\sum_{i=28}^{54}x_i\u_i\u_i\T$; then $(\Mb_{23})_{rs}\ge24-18=6>0$ for any $(r,s)\in\mathcal{K}_{24}$.
But for all $(r,s)\in\mathcal{K}_{24}$ there is exactly one row $\u\T$ of
$\left[\begin{array}{c}\! \Ub_2\! \\ \! \Ub_3\! \\ \end{array}\right]$
 satisfying $u_ru_s>0$ (which then must be a row of $\Ub_2$). This row $\u\T$ also satisfies $u_\rho u_\sigma>0$ for exactly one other $(\rho,\sigma)\in\mathcal{K}_{24}\setminus \lk (r,s)\rk$, e.g., $\u = \e_1+\e_5+\e_{10}$ for $\lk (1,10), (5,10)\rk \subset\mathcal{K}_{24}$. We thus conclude that $x_i\ge6$ must hold for all $i\in[55:63]$. However, if now $x_i=0$ for all $i\in [1:27]$, we derive $x_i=6$ for all $i\in [28:54]$ from the considerations on $\mathcal{K}_{6}$ and hence $(\Mb_{23})_{rs} = 24-18=6$ with equality in this case, which in turn implies $x_i=6$ for all $i\in [55:63]$. But then $x_{64}=19-3\cdot6=1>0$ must hold. Indeed, for any $(r,s)\in\mathcal{K}_{19}$, there are exactly three rows $\u$ of $\Ub_2$ such that $u_ru_s>0$ (in case $(r,s)=(1,5)$, these are $(\e_1+\e_5+\e_{10})\T$, $(\e_1+\e_5+\e_{11})\T$ and $(\e_1+\e_5+\e_{12})\T$), and obviously, no row $\u\T$ of $\left[\begin{array}{c}\! \Ub_0\! \\ \! \Ub_1\! \end{array}\right]$ can satisfy $u_ru_s>0$. Summarizing, we have $x_i>0$ for all $i\in[28:63]$, and $x_i>0$ for at least one $i\in[1:27]\cup\{64\}$, which means $\cpr \Mb\ge37$.
From the definition of $\Mb$ we see that $\cpr \Mb\le37$, so we finally conclude $\cpr \Mb=37$.
\ws_cor_1{Figure~2 depicts a graphical summary of the preceding arguments to be read from left to right. The arrows indicate that e.g.\ exactly 3 rows in $\Ub_1$ contribute to any element of $\mathcal{K}_{24}$, and any row in $\Ub_1$ contributes to exactly two elements of $\mathcal{K}_{24}$ (we say that a row $\u_i$ of $\Ub$ {\em contributes} to $(r,s)\in\mathcal{K}_m$ iff $(\u_i\u_i\T)_{rs}>0$ and in this case we call $x_i$ the contribution of $\u_i$ to $(r,s)$; \corrr{the position of arrows is meant to be symbolic and refers to black dots only which constitute $\mathcal{K}_m$}). The circles contain bounds on contributions, e.g.~the left semicircle on top of $\mathcal{K}_{24}$ tells us that the total contribution of the rows in $\Ub_1$ to any element of $\mathcal{K}_{24}$ is positive and at most 18.} \\[.5cm]

%{\centering
%\includegraphics[scale=.88]{roadmap1_djl.pdf}}

\begin{figure}[h]
\centering
\ru_cor_1{
\includegraphics[scale=.88]{roadmap.eps}%{roadmap1_djl.pdf}
\caption{Graphical summary of the explicit construction; see text.}
}
 %The arrows indicate that e.g.\ exactly 3 rows in $\Ub_1$ contribute to any element of %$\mathcal{K}_{24}$, and any row in $\Ub_1$ contributes to exactly two elements of $\mathcal{K}_{24}$.  %The left semicircle on top of $\mathcal{K}_{24}$ tells us that the %total contribution of the rows in $\Ub_1$ to any element of $\mathcal{K}_{24}$ is positive and at most 18.}
 \label{fig:roadmap}
\end{figure}

\end{appendix}

\corrr{{\bf Acknowledgement.} The authors are indebted to two anonymous referees for their valuable remarks which helped us a lot to improve the presentation of our paper.}

%\newpage

\begin{thebibliography}{99}
\small

\bibitem{Berm98}
Abraham Berman and Naomi Shaked-Monderer.
\newblock Remarks on completely positive matrices.
\newblock {\em Linear and Multilinear Algebra}, 44:149--163, 1998.

\bibitem{Berm03}
Abraham Berman and Naomi Shaked-Monderer.
\newblock {\em Completely positive matrices}.
\newblock World Scientific Publishing Co. Inc., River Edge, NJ, 2003.

\bibitem{Berr07a}
Michael~{W.} Berry, Murray Browne, Amy~{N.} Langville, {V.}~Paul Pauca, and Robert~J. Plemmons.
\newblock Algorithms and applications for approximate nonnegative matrix factorization.
\newblock {\em Comput. Stat. \& Data Anal.}, 52(1):155�--173, 2007.

\bibitem{BSU} Immanuel M. Bomze, Werner Schachinger and Reinhard Ullrich.
\newblock From seven to eleven: completely positive matrices with high cp-rank.
\newblock \ws_cor_1{{\em Linear Algebra Appl.}, 459:208--221, 2014.}
%Also available at:
%\newblock http://www.optimization-online.org/DB\_HTML/2014/01/4206.html, 2014.

\bibitem{Bomz11a}
Immanuel~M. Bomze.
\newblock Copositive optimization -- recent developments and applications.
\newblock {\em European J. Oper. Res.}, 216:509--520, 2012.

\bibitem{Bomz11b}
Immanuel~M. Bomze, Werner Schachinger, and Gabriele Uchida.
\newblock Think co(mpletely~)positive~! -- matrix properties, examples and a
  clustered bibliography on copositive optimization.
\newblock {\em J. Global Optim.}, 52:423--445, 2012.


\bibitem{Bure11}
Samuel Burer.
\newblock Copositive programming.
\newblock In Miguel~F. Anjos and Jean~Bernard Lasserre, editors, {\em Handbook
  of Semidefinite, Cone and Polynomial Optimization: Theory, Algorithms,
  Software and Applications}, International Series in Operations Research and
  Management Science, pages 201--218. Springer, New York, 2012.

\bibitem{Drew96a}
John~H. Drew and Charles~R. Johnson.
\newblock The no long odd cycle theorem for completely positive matrices.
\newblock In {\em Random discrete structures}, volume~76 of {\em IMA Vol. Math.
  Appl.}, pages 103--115. 1996.

\bibitem{Drew94}
John~H. Drew, Charles~R. Johnson, and Raphael Loewy.
\newblock Completely positive matrices associated with {$M$}-matrices.
\newblock {\em Linear Multilinear Algebra}, 37(4):303--310, 1994.

\bibitem{Duer10}
Mirjam D{\"u}r.
\newblock Copositive programming --- a survey.
\newblock In Moritz Diehl, Francois Glineur, Elias Jarlebring, and Wim
  Michiels, editors, {\em Recent Advances in Optimization and its Applications
  in Engineering}, pages 3--20. Springer, Berlin Heidelberg New York, 2010.


\bibitem{H95}
Paul R. Halmos. \emph{Linear algebra problem book}, The Dolciani Mathematical Expositions, 16. Mathematical Association of America, Washington, DC, 1995.

\bibitem{Hild14a}
Roland Hildebrand.
\newblock Minimal zeros of copositive matrices.
\newblock \corrr{{\em Linear Algebra Appl.}, 459: 154-�174, 2014.}
%Preprint, {\tt http://arxiv.org/abs/1401.0134}, 2014.

\bibitem{KhatriRaoRef} Shuangzhe Liu.
\newblock Matrix results on the {K}hatri-{R}ao and {T}racy-{S}ingh products.
\newblock {\em Linear Algebra Appl.}, 289:267--277, 1999.

\bibitem{Loew03}
Raphael Loewy and Bit-Shun Tam.
\newblock C{P} rank of completely positive matrices of order~5.
\newblock {\em Linear Algebra Appl.}, 363:161--176, 2003.

\bibitem{Shak13b}
Naomi Shaked-Monderer, Abraham Berman, Immanuel~M. Bomze, Florian Jarre, and
  Werner Schachinger.
\newblock New results on the cp rank and related properties of
  co(mpletely~)positive matrices.
\newblock {\em Linear Multilinear Algebra}, to appear. Also available
  at:~arxiv.org/abs/1305.0737, 2013.

\bibitem{Shak13a}
Naomi Shaked-Monderer, Immanuel~M. Bomze, Florian Jarre, and Werner
  Schachinger.
\newblock On the cp-rank and minimal cp factorizations of a completely positive
  matrix.
\newblock {\em SIAM J. Matrix Anal. Appl.}, 34(2):355--�368, 2013.

\bibitem{Vava09a}
Stephen~{A.} Vavasis.
\newblock On the complexity of nonnegative matrix factorization.
\newblock {\em SIAM J. Optim.}, 20(3):1364�--1377, 2009.

\end{thebibliography}


\end{document}





